\section{The threshold one half, two-dimensional models and limits of methods}
\label{app:obstructions}

This appendix proves the statements about the hypothesis of
Section~\ref{sec:massgen} and about other approaches to (D) that are used in
the text. Section~\ref{ssec:gausschain} shows that the threshold $c<\frac12$ of
Theorem~\ref{thm:D} is the exact limit of the slab mechanism, the boundary value
being attained by a massless Gaussian chain. Sections~\ref{ssec:twodfat}
and~\ref{ssec:strips} treat two dimensions: the slab hypothesis for fat boxes
holds along the weak-coupling tuning, uniformly in the boundary condition, while
the same condition for strips of width one fails. Sections~\ref{ssec:curvlimits}
to~\ref{ssec:transfer} prove that four other routes cannot give (D) without an
input of the same strength: curvature of the Wilson measure, entropy
factorisation, tensorisation over blocks of fixed lattice size, and the
transfer matrix taken alone.

\subsection{The Gaussian chain}\label{ssec:gausschain}

For a probability measure and $\sigma$-algebras $\mathcal F_1,\mathcal F_2$ let
$\rho(\mathcal F_1,\mathcal F_2)=\sup|E[XY]|$, the supremum over centred $X\in
L^2(\mathcal F_1)$ and $Y\in L^2(\mathcal F_2)$ with $E[X^2]=E[Y^2]=1$; this is
the maximal correlation of Definition~\ref{def:rho}, without the constraint of
gauge invariance.

\begin{lemma}\label{lem:gmarkov}
Let $(U,V,M,W,X)$ be random variables with values in standard Borel spaces whose
joint law has the form $\psi_1(u,v)\psi_2(v,m,w)\psi_3(w,x)\,d\lambda$ for a
product measure $\lambda$ and nonnegative functions $\psi_i$. Then
$\rho(\sigma(U,V),\sigma(W,X))\le\rho(\sigma(V),\sigma(W))$.
\end{lemma}

\begin{proof}
By the factorisation, the conditional law of $U$ given $(V,M,W,X)$ has density
proportional to $\psi_1(\cdot,V)$ and that of $X$ given $(U,V,M,W)$ has density
proportional to $\psi_3(W,\cdot)$, and given $(V,M,W)$ the variables $U$ and
$X$ are conditionally independent. Let $f=f(U,V)$ and $g=g(W,X)$ be centred and
square integrable. Then $E[f\mid V,M,W,X]=\tilde f(V)$ with $\tilde f=E[f\mid
V]$, and $E[g\mid U,V,M,W]=\tilde g(W)$ with $\tilde g=E[g\mid W]$. Hence
$E[fg]=E[E[fg\mid V,M,W]]=E[\tilde f(V)\tilde g(W)]$, where $\tilde f$ and
$\tilde g$ are centred with $\|\tilde f\|\le\|f\|$ and $\|\tilde g\|\le\|g\|$.
\end{proof}

\begin{lemma}[Gebelein]\label{lem:gebelein}
If $(V,W)$ is a nondegenerate Gaussian vector in $\R^2$ with correlation
coefficient $r$, then $\rho(\sigma(V),\sigma(W))=|r|$.
\end{lemma}

\begin{proof}
The choice $X=V$, $Y=\pm W$, standardised, gives $\rho\ge|r|$. For the converse
standardise $V$ and $W$. Let $H_n$ be the Hermite polynomials with generating
function $\sum_nH_n(x)s^n/n!=e^{sx-s^2/2}$; then $h_n=H_n/\sqrt{n!}$ is a
complete orthonormal system of $L^2(\R,\gamma)$, $\gamma$ the standard Gaussian
law~\cite[Sec.~1.1.1]{Nua06}. From $E[e^{sV-s^2/2}e^{uW-u^2/2}]=e^{rsu}$ and a
comparison of coefficients one gets $E[h_j(V)h_k(W)]=\delta_{jk}r^j$; the double
series converges absolutely, since $H_n(x)=E(x+iZ)^n$ with $Z\sim\gamma$ gives
$|H_n(x)|\le E(|x|+|Z|)^n$ and $Ee^{t(|V|+|W|+2|Z|)}<\infty$ for every $t$.
For centred $f=\sum_{j\ge1}a_jh_j$ and $g=\sum_{j\ge1}b_jh_j$ in $L^2(\gamma)$
the form $(f,g)\mapsto E[f(V)g(W)]$ is bounded by $\|f\|\|g\|$, so passing to
the limit in the partial sums, $E[f(V)g(W)]=\sum_{j\ge1}a_jb_jr^j$, whose
modulus is at most $|r|\sum_j|a_jb_j|\le|r|\|f\|\|g\|$.
\end{proof}

The inequality $\rho\le|r|$ is Gebelein's theorem~\cite{Geb41}.

\emph{The chain.} Let $\mu\ge0$ and define $\kappa\ge0$ by $2\cosh\kappa=2+
\mu^2$. For integers $p<q$ with $q-p\ge2$ and real boundary values $z_p,z_q$ let
$\nu$ be the Gaussian law of $(z_{p+1},\dots,z_{q-1})$ with density proportional
to
\[
\exp\Bigl(-\tfrac12\sum_{j=p}^{q-1}(z_{j+1}-z_j)^2-\tfrac12\mu^2\sum_{j=p+1}^{q-1}
z_j^2\Bigr),
\]
a nearest-neighbour Gaussian chain, massless for $\mu=0$. Its precision matrix
on the $q-p-1$ interior sites is $\mathsf L=\operatorname{tridiag}(-1,2+\mu^2,
-1)$, independent of $z_p$ and $z_q$, which only shift the mean, and its
conditional laws on subintervals are chains of the same form.

\begin{lemma}\label{lem:green}
Put $T=q-p$ and, for $0\le i\le T$, $u(i)=\sinh(\kappa i)$ and $v(i)=\sinh(
\kappa(T-i))$ if $\mu>0$, $u(i)=i$ and $v(i)=T-i$ if $\mu=0$. For $0<i\le j<T$,
$\Cov_\nu(z_{p+i},z_{p+j})=A\,u(i)v(j)$ with $A>0$ independent of $i$ and $j$;
for $\mu=0$, $A=1/T$. Hence the correlation of $z_{p+i}$ and $z_{p+j}$ is
$(u(i)v(j)/(v(i)u(j)))^{1/2}$.
\end{lemma}

\begin{proof}
Both $u$ and $v$ satisfy $(2+\mu^2)f(l)=f(l-1)+f(l+1)$ (for $\mu>0$ because
$\sinh(x-\kappa)+\sinh(x+\kappa)=2\cosh\kappa\sinh x$), $u(0)=v(T)=0$, and both
are positive on $(0,T)$. The covariance matrix is $G=\mathsf L^{-1}$, which is
symmetric. For fixed $j$ the column $G(\cdot,j)$, extended by $0$ at $0$ and
$T$, satisfies the recursion at every $l\ne j$, so $G(i,j)=\alpha_ju(i)$ for
$i\le j$ and $G(i,j)=\beta_jv(i)$ for $i\ge j$: a solution of the second-order
recursion on $[0,j]$ that vanishes at $0$ is a multiple of $u$, and similarly
for $v$. For $i<j$, $\alpha_ju(i)=G(i,j)=G(j,i)=\beta_iv(j)$, so $\alpha_j/v(j)
=\beta_i/u(i)$ for all $i<j$, a constant $A$; and $\alpha_ju(j)=\beta_jv(j)$
gives $G(j,j)=Au(j)v(j)$, so $A>0$. For $\mu=0$ the $j$-th entry of $\mathsf L
G(\cdot,j)$ equals one, that is $A[2j(T-j)-(j-1)(T-j)-j(T-j-1)]=AT=1$.
\end{proof}

\emph{Slab geometry.} Let $s\ge1$, $s\le w\le2s-1$ and $M\ge1$. Take the sites
$I=[0,sM+w)$ with the two exterior sites $-1$ and $sM+w$, the blocks
$R_t=[st,st+w)$ for $0\le t\le M$, and, for $t<M$, $Q_t=R_t\cup R_{t+1}$ split
into $Q_t^-=R_t\setminus R_{t+1}=[st,st+s)$, $Q_t^0=R_t\cap R_{t+1}=[st+s,
st+w)$ and $Q_t^+=R_{t+1}\setminus R_t=[st+w,st+s+w)$. The condition $w\le2s-1$
says that $R_t$ and $R_{t+2}$ contain no adjacent sites, so that, exactly as in
Proposition~\ref{prop:slab}, the resampling moves of non-adjacent blocks
commute; Section~\ref{sec:massgen} uses $s=d$ and $w=2d-1$. Let $\nu$ be the
chain on $I$ with any exterior values, $P_t$ the conditional expectation given
the sites outside $R_t$, and $\rho_t$ the supremum, over all values outside
$Q_t$, of $\rho(\sigma(z_{Q_t^-}),\sigma(z_{Q_t^+}))$ under the conditional law
of $z_{Q_t}$.

\begin{proposition}[the Gaussian chain]\label{prop:gausschain}
\leavevmode
\begin{enumerate}[label=\textup{(\roman*)},leftmargin=*]
\item For every $t<M$, $\rho_t=\sinh(\kappa s)/\sinh(\kappa(w+1))$ if $\mu>0$,
and $\rho_t=s/(w+1)$ if $\mu=0$. For $w=2s-1$ this is $1/(2\cosh(\kappa s))$,
and $\frac12$ for $\mu=0$. For $\mu=0$, $\rho_t\ge\frac12$ in every admissible
geometry, with equality only for $w=2s-1$.
\item For $\mu=0$ the best constant $\gamma_M$ in the inequality $\gamma_M
\Var_\nu(F)\le\sum_{t=0}^M\nu((F-P_tF)^2)$ for all $F\in L^2(\nu)$ satisfies
$\gamma_M\le16/(3M)$.
\item For a subinterval split into consecutive nonempty intervals of $A'$, $r$
and $C'$ sites, with the sites outside it conditioned, the maximal correlation
of the two outer intervals is
\[
\Bigl(\frac{\sinh(\kappa A')\sinh(\kappa C')}{\sinh(\kappa(A'+r+1))\sinh(\kappa(
C'+r+1))}\Bigr)^{1/2}\le e^{-\kappa(r+1)}\quad\text{if }\mu>0,
\]
with supremum $e^{-\kappa(r+1)}$ over $A'$ and $C'$, and $(A'C'/((A'+r+1)(C'+r
+1)))^{1/2}$ if $\mu=0$, with supremum $1$.
\end{enumerate}
\end{proposition}

\begin{proof}
(i) Given the values outside $Q_t$, the law of $z_{Q_t}$ is the chain with
$p=st-1$, $q=st+s+w$, $T=s+w+1$. Its density factorises as in
Lemma~\ref{lem:gmarkov} with $U=z_{[st,st+s-1)}$, $V=z_{st+s-1}$,
$M=z_{Q^0_t}$, $W=z_{st+w}$ and $X=z_{(st+w,st+s+w)}$, the interactions being
between nearest neighbours; so $\rho_t\le\rho(\sigma(V),\sigma(W))$, with
equality since $V$ and $W$ are functions of the two outer parts. By
Lemma~\ref{lem:gebelein} this is the modulus of the correlation of $V$ and $W$,
which by Lemma~\ref{lem:green} with $i=s$, $j=w+1$ and $T=s+w+1$ equals
$(u(s)v(w+1)/(v(s)u(w+1)))^{1/2}$. Here $v(w+1)=u(s)$ and $v(s)=u(w+1)$, which
gives the formulas; they do not depend on the conditioning values. For $w=2s-1$,
$\sinh(\kappa s)/\sinh(2\kappa s)=1/(2\cosh\kappa s)$. For $\mu=0$, $w+1\le2s$
gives $s/(w+1)\ge\frac12$, with equality only when $w=2s-1$.

(ii) Let $T=sM+w+1\ge sM\ge2$ and take $F=z_j$ with $j=\lfloor T/2\rfloor-1\in
I$. By Lemma~\ref{lem:green} with $p=-1$, $\Var_\nu(F)=i(T-i)/T$ with $i=
\lfloor T/2\rfloor$, so $\Var_\nu(F)\ge(T^2-1)/(4T)\ge3T/16$. If $j\notin R_t$,
then $P_tF=F$. If $j\in R_t$, then $\nu((F-P_tF)^2)$ is the mean conditional
variance of $z_j$ given $z_{st-1}$ and $z_{st+w}$, the variance of a chain of
$w+1$ steps at an interior site, which is at most $(w+1)/4\le s/2$ by
Lemma~\ref{lem:green}. The $t$ with $j\in R_t$ satisfy $(j-w)/s<t\le j/s$, an
interval of length $w/s<2$, so there are at most two. Hence $\sum_t\nu((F-P_t
F)^2)\le s$, and $\gamma_M\le16s/(3T)\le16/(3M)$.

(iii) Put the left conditioned site at $0$. The outer parts are $[1,A']$ and
$[A'+r+1,A'+r+C']$, and the right conditioned site is $T=A'+r+C'+1$. As in (i),
by Lemmas~\ref{lem:gmarkov} and~\ref{lem:gebelein}, the maximal correlation is
the modulus of the correlation of the inner ends $z_{A'}$ and $z_{A'+r+1}$,
which by Lemma~\ref{lem:green} is $(u(A')v(A'+r+1)/(v(A')u(A'+r+1)))^{1/2}$.
With $u_0(x)=\sinh\kappa x$, respectively $u_0(x)=x$, one has $u(i)=u_0(i)$ and
$v(i)=u_0(T-i)$, so $v(A'+r+1)=u_0(C')$ and $v(A')=u_0(C'+r+1)$, which gives the
formulas. For $\mu>0$ the inequality $\sinh x/\sinh(x+y)\le e^{-y}$ for $x,y\ge
0$, which is equivalent to $e^{-x-y}\le e^{-x+y}$, bounds both factors by
$e^{-\kappa(r+1)}$, and both tend to $e^{-\kappa(r+1)}$ as $A',C'\to\infty$.
For $\mu=0$ the expression tends to $1$.
\end{proof}

Three consequences explain the role of the chain in Section~\ref{sec:massgen}.
First, for $w=2s-1$ the massless chain satisfies the hypothesis of
Proposition~\ref{prop:slab} with $c=\frac12$ for every exterior condition, its
non-adjacent moves commute, and yet by (ii) no inequality $\gamma\Var_\nu(F)\le
\sum_t\nu((F-P_tF)^2)$ holds with $\gamma>0$ independent of $M$; in any geometry
with commuting non-adjacent moves the same holds for every $c\ge s/(w+1)\ge
\frac12$. So the threshold $\frac12$ of Lemma~\ref{lem:chain} cannot be raised
by an argument that uses only these two properties, and no slab geometry of
this kind improves it. Second, for $w=2s-1$ the slab constant of the chain is
$1/(2\cosh\kappa s)$, which is below $\frac12$ if and only if $\mu>0$: along a
tuning $s=\ell/(2a)$ with lattice mass $\mu=\mu_{\rm phys}a$ one has $\kappa/
\mu\to1$, so the constant tends to $1/(2\cosh(\mu_{\rm phys}\ell/2))<\frac12$.
In the chain the hypothesis is exactly a mass at the scale $\ell/2$. Third, the
condition $(\mathrm M^*_r(c))$ with a middle of width $r$ has in the chain
the best constant $e^{-\kappa(r+1)}$ by (iii), and for $\mu=0$ it fails for
every $c<1$ and every $r$; the iteration of Proposition~\ref{prop:alldecay} reproduces the
true decay up to the integer part of the ratio of widths, so it loses nothing,
but it only transports a mass present at the starting scale.

\subsection{Two dimensions: fat boxes}\label{ssec:twodfat}

In this section and the next, $G$ is a compact group with normalised Haar
measure $\lambda$, also on products of copies of $G$, and $\varphi$ is a
positive continuous class function on $G$ with $\varphi(g^{-1})=\varphi(g)$; for
$G=\SU(2)$ the last condition holds for every class function, since every
element of $\SU(2)$ is conjugate to its inverse. Multiplying $\varphi$ by a
constant changes no measure, so we normalise $\int\varphi\,d\lambda=1$. For
class functions put $(\alpha*\beta)(g)=\int\alpha(gy^{-1})\beta(y)\,dy$ and let
$\alpha^{*k}$ be the $k$-fold convolution power.

Links of $\Z^2$ are $(y,\nu)$, $\nu\in\{1,2\}$, from $y$ to $y+e_\nu$, and the
plaquette with base point $z$ has holonomy $U_p=U_{(z,1)}U_{(z+e_1,2)}
U_{(z+e_2,1)}^{-1}U_{(z,2)}^{-1}$. Boxes $Q=([\alpha_1,\alpha_1+s_1)\times[
\alpha_2,\alpha_2+s_2))\cap\Z^2$, the sets $V(Q)$, the measures $\nu^\omega_Q$,
the groups $\mathcal G_Q$, the spaces $\mathcal A^\omega_Q$, splittings and the
maximal correlation $\rho^\omega(Q_-,Q_+)$ are defined as in
Section~\ref{sec:massgen} and Definition~\ref{def:rho}, with $\Z^2$ in place of
$\Z^4$; a vertex $v$ is interior to $Q$ when its four incident links lie in
$V(Q)$, that is, when $v$, $v-e_1$ and $v-e_2$ lie in $Q$. The
\emph{two-dimensional form} of $(\mathrm M_{d,m}(c))$ is the condition
$\rho^\omega(Q_-,Q_+)\le c$ for every box of $\Z^2$ whose two widths lie in
$[2d-1,md-1]$, every coordinate in which it has width $3d-1$, the splitting into
widths $d$, $d-1$, $d$ in that coordinate, and every exterior configuration
$\omega$.

\begin{lemma}[merging]\label{lem:merge}
Let $\alpha,\beta$ be continuous class functions on $G$.
\begin{enumerate}[label=\textup{(\roman*)},leftmargin=*]
\item $\alpha*\beta$ is a continuous class function, positive if $\alpha$ and
$\beta$ are, and $\int\alpha(Ax)\beta(x^{-1}B)\,dx=(\alpha*\beta)(AB)$ for all
$A,B\in G$.
\item Let $\alpha_0,\dots,\alpha_M$ be continuous class functions, $M\ge1$, and
$g_j,c_j,V_0,V_{M+1}\in G$. Then
\[
\int\prod_{j=0}^M\alpha_j\bigl(V_{j+1}g_jV_j^{-1}c_j\bigr)\,dV_1\cdots dV_M=
(\alpha_0*\cdots*\alpha_M)\bigl(V_{M+1}\,g_M\cdots g_0\,V_0^{-1}\,c_0\cdots c_M
\bigr).
\]
\end{enumerate}
\end{lemma}

\begin{proof}
(i) The substitution $y=x^{-1}B$ preserves Haar measure and turns the integral
into $\int\alpha(ABy^{-1})\beta(y)\,dy$. For $h\in G$ the substitution $y=hy'
h^{-1}$ gives $(\alpha*\beta)(hgh^{-1})=\int\alpha(hgy'^{-1}h^{-1})\beta(hy'
h^{-1})\,dy'=(\alpha*\beta)(g)$. Continuity and positivity are clear.

(ii) For $M=1$: since $\alpha_0$ and $\alpha_1$ are class functions,
$\alpha_0(V_1g_0V_0^{-1}c_0)=\alpha_0(g_0V_0^{-1}c_0V_1)$ and $\alpha_1(V_2g_1
V_1^{-1}c_1)=\alpha_1(V_1^{-1}c_1V_2g_1)$, so by (i) the integral over $V_1$ is
$(\alpha_0*\alpha_1)(g_0V_0^{-1}c_0c_1V_2g_1)=(\alpha_0*\alpha_1)(V_2g_1g_0
V_0^{-1}c_0c_1)$. For $M\ge2$ integrate $V_1,\dots,V_{M-1}$ first; by
induction the result is $\beta(V_Mg_{M-1}\cdots g_0V_0^{-1}c_0\cdots c_{M-1})$
with $\beta=\alpha_0*\cdots*\alpha_{M-1}$, and the case $M=1$ with
$g_{M-1}\cdots g_0$ and $c_0\cdots c_{M-1}$ in place of $g_0$ and $c_0$ gives
the claim.
\end{proof}

Fix a box, after translation $Q=[0,s_1)\times[0,s)$ with $s_1=3d-1$, $d\ge2$,
$s\ge1$, split in the first coordinate into $Q_-=[0,d)\times[0,s)$, $Q_0=
[d,2d-1)\times[0,s)$ and $Q_+=[2d-1,3d-1)\times[0,s)$. Write $h(x,y)=
\sigma_{((x,y),1)}$ and $v(x,y)=\sigma_{((x,y),2)}$. The interior vertices are
$I=[1,s_1-1]\times[1,s-1]$. Let
\[
F'=\{((x,y),2):1\le x\le s_1-1,\ 0\le y\le s-2\},
\]
the vertical links entering the interior vertices from below ($F'=\emptyset$
if $s=1$), and $E'=V(Q)\setminus F'$: all $h(x,y)$, all $v(0,y)$, and
$v(x,s-1)$ for $x\ge1$.

\begin{lemma}[axial gauge]\label{lem:axial}
For $\sigma\in G^{V(Q)}$ define $H(\sigma)\in\mathcal G_Q$ recursively in the
second coordinate by $H_u=\sigma_{(u-e_2,2)}^{-1}H_{u-e_2}$ for $u\in I$, with
$H_w=1$ for $w\notin I$, and put $\tau(\sigma)=(H(\sigma)^{-1}\sigma)|_{E'}$.
\begin{enumerate}[label=\textup{(\roman*)},leftmargin=*]
\item $(H(\sigma)^{-1}\sigma)_e=1$ for $e\in F'$, and $H(\sigma)$ depends only
on $\sigma_{F'}$.
\item Under $\lambda$ on $G^{V(Q)}$, $\tau(\sigma)$ has law $\lambda$ on $G^{E'}$.
\item Under $\nu=\nu^\omega_Q$, $\tau(\sigma)$ has density proportional to
$\mathcal W(\tau)=\prod_{p\cap V(Q)\ne\emptyset}\varphi((\tau,1_{F'},\omega)_p)$.
\item If $X\in\mathcal A^\omega_Q$ is measurable with respect to
$\sigma_{V(Q_-)}$, there is a Borel function $\bar X$ with $X=\bar X(\tau(
\sigma)_{E'\cap V(Q_-)})$ $\nu$-almost surely; likewise for $Q_+$.
\end{enumerate}
\end{lemma}

\begin{proof}
(i) For $u\in I$, $(H^{-1}\sigma)_{(u-e_2,2)}=H_{u-e_2}^{-1}\sigma_{(u-e_2,2)}
H_u=1$ by the definition of $H_u$. The recursion terminates, since $u-e_2$ has
a smaller second coordinate, and it uses only the links $(u-je_2,2)$, $1\le j
\le u_2$, which lie in $F'$.

(ii) Under $\lambda$ the families $\sigma_{F'}$ and $\sigma_{E'}$ are
independent with laws $\lambda$. Given $\sigma_{F'}$, $H$ is fixed, and for a
link $e$ of $E'$ from $x$ to $y$, $\tau_e=H_x^{-1}\sigma_eH_y$ is a two-sided
translate of $\sigma_e$. Translations preserve $\lambda$, so the conditional
law of $\tau$ given $\sigma_{F'}$ is $\lambda$, and so is its law.

(iii) Let $w(\sigma)=\prod_{p\cap V(Q)\ne\emptyset}\varphi((\sigma,\omega)_p)$.
Every $h\in\mathcal G_Q$ fixes the links outside $V(Q)$, whose end points are
not interior, so $(h\sigma,\omega)=h\cdot(\sigma,\omega)$ and each plaquette
holonomy is conjugated; as $\varphi$ is a class function, $w(h\sigma)=
w(\sigma)$. Hence $w(\sigma)=w(H(\sigma)^{-1}\sigma)=\mathcal W(\tau(\sigma))$
by (i), and for bounded Borel $\Psi$, by (ii), $\int\Psi(\tau(\sigma))
w(\sigma)\,d\lambda(\sigma)=\int\Psi(\tau)\mathcal W(\tau)\,d\lambda(\tau)$.

(iv) The group $\mathcal G_Q\cong G^I$ acts on $G^{V(Q_-)}$ link by link. Let
$X=f(\sigma_{V(Q_-)})$ almost surely with $f$ Borel, and put $\bar f(\eta)=
\int f(h\eta)\,dh$ where the integral converges absolutely and $\bar f=0$
elsewhere; by invariance of Haar measure both the set of convergence and $\bar f$
are invariant under every $k\in\mathcal G_Q$. Since $\nu$ and $\lambda$ have the
same null sets, the density being bounded above and below, and $f(h\eta)=
f(\eta)$ for $\nu$-almost every configuration for each $h$, Fubini's theorem
gives $\bar f=f$ almost surely. Then $X=\bar f(\sigma_{V(Q_-)})=\bar f((H(
\sigma)^{-1}\sigma)_{V(Q_-)})$ almost surely, and by (i) the right side is a
function of $\tau(\sigma)_{E'\cap V(Q_-)}$, the coordinates in $F'$ being $1$.
\end{proof}

\begin{lemma}[near-product densities]\label{lem:nearproduct}
Let $(S_1,\mu_1)$ and $(S_2,\mu_2)$ be probability spaces and $\pi$ a
probability measure on $S_1\times S_2$ with density $r$ with respect to
$\mu_1\otimes\mu_2$, $|r-1|\le\delta<1$. If $X\in L^2(\pi)$ depends only on
$s_1$, $Y\in L^2(\pi)$ only on $s_2$, and $\pi(X)=\pi(Y)=0$, then $|\pi(XY)|\le
\frac{\delta+\delta^2}{1-\delta}\|X\|_{L^2(\pi)}\|Y\|_{L^2(\pi)}$. If $r=k/Z'$
with $1-\varepsilon\le k\le1+\varepsilon$ and $Z'=\int k\,d(\mu_1\otimes\mu_2)$,
then $\delta=2\varepsilon/(1-\varepsilon)$ is admissible, and for $\varepsilon\le
\frac1{10}$ the bound is at most $4\varepsilon$.
\end{lemma}

\begin{proof}
The marginal density $m_1(s_1)=\int r(s_1,\cdot)\,d\mu_2$ lies in $[1-\delta,1+
\delta]$, so $\int X^2d\mu_1\le(1-\delta)^{-1}\|X\|^2_{L^2(\pi)}$; and $\int X\,
d\mu_1=\int X(1-m_1)\,d\mu_1$ since $\int Xm_1\,d\mu_1=\pi(X)=0$, so $|\int X\,
d\mu_1|\le\delta\int|X|\,d\mu_1$. The same holds for $Y$. Now $\pi(XY)=\iint XY
(r-1)\,d\mu_1d\mu_2+\int X\,d\mu_1\int Y\,d\mu_2$, and by the Cauchy--Schwarz
inequality the two terms together are at most $(\delta+\delta^2)(\int X^2d\mu_1
\int Y^2d\mu_2)^{1/2}$, which gives the first claim. If $r=k/Z'$, then $Z'\in
[1-\varepsilon,1+\varepsilon]$ and $r\in[\frac{1-\varepsilon}{1+\varepsilon},
\frac{1+\varepsilon}{1-\varepsilon}]$, so $|r-1|\le2\varepsilon/(1-\varepsilon)$.
For $\varepsilon\le\frac1{10}$, $\delta\le\frac{20}9\varepsilon\le\frac29$ and
$\frac{\delta(1+\delta)}{1-\delta}\le\frac{11}7\delta\le\frac{220}{63}
\varepsilon<4\varepsilon$.
\end{proof}

For $G=\SU(2)$ write $\chi_j$ for the character of spin $j\in\frac12\Z_{\ge0}$,
$d_j=2j+1$ and $c_j=j(j+1)$, and let
\[
p_t=\sum_jd_je^{-tc_j}\chi_j\qquad(t>0)
\]
be the heat kernel. It is a uniformly convergent series of real class
functions, $\int p_t=1$, $p_t*p_{t'}=p_{t+t'}$ because $\chi_i*\chi_j=
\delta_{ij}\chi_j/d_j$, and $p_t>0$, as the heat kernel of a compact connected
Riemannian manifold~\cite[Chs.~7--8]{Gri09}. It is the weight of the
\emph{heat-kernel action}, for which the lattice theory in two dimensions
reproduces the continuum Yang--Mills measure on every lattice
\cite{GKS89,Dri89,Lev03}, and the two-dimensional tuning is $t=g^2a^2$ with a
coupling $g$ of the dimension of a mass. With $n=2j+1$, $c_j=(n^2-1)/4$ and
$|\chi_j|\le d_j$,
\begin{equation}\label{eq:hk}
\|p_\theta-1\|_\infty\le\sum_{n\ge2}n^2e^{-\theta(n^2-1)/4}\le\sum_{k\ge1}
(k+1)^2q^k\le\frac{4q}{(1-q)^3},\qquad q=e^{-3\theta/4},
\end{equation}
using $n^2-1\ge3(n-1)$ and $\sum_{k\ge1}(k+1)^2q^k=(4q-3q^2+q^3)/(1-q)^3$. For
$\theta\ge6$, $q\le e^{-4.5}<\frac1{20}$ and the right side is at most $5q\le
0.06$.

\begin{theorem}[two dimensions, fat boxes]\label{thm:twod}
\leavevmode
\begin{enumerate}[label=\textup{(\roman*)},leftmargin=*]
\item Let $d\ge2$, $s\ge1$, let $Q$ be a box of $\Z^2$ of widths $3d-1$ and $s$,
in either order, split in the coordinate of width $3d-1$ into widths $d$,
$d-1$, $d$, and put $N=1+(d-1)(s+1)$ and $\varepsilon_N=\|\varphi^{*N}-1
\|_\infty$. If $\varepsilon_N\le\frac1{10}$, then $\rho^\omega(Q_-,Q_+)\le4
\varepsilon_N$ for every exterior configuration $\omega$.
\item Let $G=\SU(2)$, $\varphi=p_t$ with $t=g^2a^2$, $d\ge2$ and $gda\ge\sqrt6$.
For every box of widths $3d-1$ and $s\ge2d-1$ and every $\omega$,
$\rho^\omega(Q_-,Q_+)\le20e^{-3Nt/4}$ with $N=1+(d-1)(s+1)$. In particular the
two-dimensional form of $(\mathrm M_{d,m}(c))$ holds for every $m\ge4$ with
\[
c=20\,e^{-3g^2(da)^2/4}.
\]
\end{enumerate}
\end{theorem}

\begin{proof}
(i) \emph{Reduction to the split in the first coordinate.} The exchange of the
two coordinates maps links to links, boxes to boxes, interior vertices to
interior vertices and $\mathcal G_Q$ to $\mathcal G_{Q'}$, and it maps the
holonomy of the plaquette at $z$ to the inverse of the holonomy of the
plaquette at the image of $z$; since $\varphi(g^{-1})=\varphi(g)$, it maps
$\nu^\omega_Q$ to $\nu^{\omega'}_{Q'}$. So we may take $Q$ as above.

\emph{Reduction to the axial model.} By Lemma~\ref{lem:axial}(iii) and (iv),
$\rho^\omega(Q_-,Q_+)$ is at most the supremum of $|\pi(\bar X\bar Y)|$ over
centred unit vectors $\bar X\in L^2(\pi)$ depending on $\tau_-=\tau_{E'\cap
V(Q_-)}$ and $\bar Y$ depending on $\tau_+=\tau_{E'\cap V(Q_+)}$, where $\pi$ is
the law on $G^{E'}$ with density proportional to $\mathcal W$.

\emph{Which plaquettes couple the three parts.} Every plaquette meeting $V(Q)$
contains a link of $E'$: the plaquette at $z\in Q$ contains $h(z)$, the
plaquette at $(-1,y)$ contains $v(0,y)$, and the plaquette at $(x,-1)$ contains
$h(x,0)$. The plaquette at $(x,y)$ contains $h(x,y)$, $h(x,y+1)$ and $v(x,y)$,
which lie in the part containing the column $x$, and $v(x+1,y)$, which lies in
the part containing the column $x+1$ and belongs to $F'$, so is fixed to $1$,
unless $y=s-1$. Hence the only plaquettes that contain free links of two
different parts are the two seam plaquettes at $(d-1,s-1)$, joining $Q_-$ and
$Q_0$, and at $(2d-2,s-1)$, joining $Q_0$ and $Q_+$, and no plaquette contains
free links of both $Q_-$ and $Q_+$. So $\mathcal W=\mathcal W_-(\tau_-)\mathcal
W_+(\tau_+)\widehat{\mathcal W}$, where $\widehat{\mathcal W}$ is the product of
the factors of the plaquettes that contain a free link of $Q_0$.

\emph{Integration of the middle.} Write $H_x=h(x,s-1)$, $V_x=v(x,s-1)$ and
$e_x=\sigma_{((x,s),1)}$, an exterior link. For a column $x\in[d,2d-2]$ the
factors of $\widehat{\mathcal W}$ are: the bottom plaquette at $(x,-1)$, which
equals $\varphi(h(x,0)^{-1}c_x)$ with $c_x$ a product of exterior links, by the
cyclic invariance of $\varphi$; the plaquettes at $(x,y)$, $0\le y\le s-2$,
which equal $\varphi(h(x,y)h(x,y+1)^{-1})$, both vertical links being in $F'$;
and the top plaquette $\varphi(H_xV_{x+1}e_x^{-1}V_x^{-1})$. Integrating
$h(x,0),\dots,h(x,s-2)$ successively with Lemma~\ref{lem:merge}(i), each step
having the form $\int\varphi(h_{y+1}^{-1}h_y)\psi(h_y^{-1}c_x)\,dh_y$, gives
$\varphi^{*s}(H_x^{-1}c_x)$, and then $\int\varphi(V_{x+1}e_x^{-1}V_x^{-1}H_x)
\varphi^{*s}(H_x^{-1}c_x)\,dH_x=\varphi^{*(s+1)}(V_{x+1}e_x^{-1}V_x^{-1}c_x)$,
again by Lemma~\ref{lem:merge}(i). The seam plaquette at $(d-1,s-1)$ is
$\varphi(H_{d-1}V_de_{d-1}^{-1}V_{d-1}^{-1})=\varphi(V_de_{d-1}^{-1}V_{d-1}^{-1}
H_{d-1})$, with $H_{d-1}$ and $V_{d-1}$ free links of $Q_-$, and the seam
plaquette at $(2d-2,s-1)$ is the top plaquette of the column $2d-2$, which
contains $V_{2d-1}$, a free link of $Q_+$. Lemma~\ref{lem:merge}(ii) with
$M=d-1$, $\alpha_0=\varphi$, $\alpha_j=\varphi^{*(s+1)}$ for $j\ge1$,
$V_0=V_{d-1}$, $c_0=H_{d-1}$, $g_j=e^{-1}_{d-1+j}$ and $c_j=c_{d-1+j}$ gives,
after a cyclic rearrangement,
\[
\int\widehat{\mathcal W}\,d\lambda(\tau_0)=\varphi^{*N}(\Phi_-\Phi_+),\qquad
\Phi_-=V_{d-1}^{-1}H_{d-1}C,\quad\Phi_+=V_{2d-1}\Gamma,
\]
with $C=c_d\cdots c_{2d-2}$ and $\Gamma=e_{2d-2}^{-1}\cdots e_{d-1}^{-1}$
exterior, and $N=1+(d-1)(s+1)$.

\emph{Conclusion.} Under $\pi$ the law of the pair $(\tau_-,\tau_+)$ has a
density proportional to
\[
\mathcal W_-(\tau_-)\,\mathcal W_+(\tau_+)\,\varphi^{*N}(\Phi_-\Phi_+).
\]
Let $\mu_1$ and $\mu_2$ be the probability measures proportional to
$\mathcal W_-\lambda$ and to $\mathcal W_+\lambda$. Then the law has density
$k/Z'$ with respect to $\mu_1\otimes\mu_2$, where
$k=\varphi^{*N}(\Phi_-\Phi_+)$ takes values in
$[1-\varepsilon_N,1+\varepsilon_N]$, and Lemma~\ref{lem:nearproduct} gives the
bound $4\varepsilon_N$.

(ii) For $s\ge2d-1$, $N\ge1+2d(d-1)\ge d^2$, so $Nt\ge g^2d^2a^2\ge6$. Since
$\varphi^{*N}=p_{Nt}$, \eqref{eq:hk} gives $\varepsilon_N\le5e^{-3Nt/4}\le0.06$,
and (i) gives $\rho^\omega\le20e^{-3Nt/4}\le20e^{-3g^2d^2a^2/4}$.
\end{proof}

The heat-kernel weight has string tension $\frac34g^2$, and $Nt$ is $g^2$ times
the area of the middle slab up to boundary terms, so the bound reads
$\rho\le20e^{-\frac34g^2\cdot\mathrm{area}}$: in two dimensions the slab
decorrelation of Section~\ref{sec:massgen} is the confinement of the middle
slab, and it holds uniformly in the boundary condition with $c\to0$ as the
physical size $da$ of the box grows. The proof rests on one fact. After the
axial gauge the plaquettes straddling each seam contain a single free link, so
the middle slab integrates out exactly to a function $\varphi^{*N}(\Phi_-\Phi_+)$
of one group element, uniformly close to $1$ once the middle slab is confined.
In four dimensions the same gauge leaves on each seam the free links of a
three-dimensional layer, and integrating out the middle produces the
$(d-1)$-step transfer operator of the box, a kernel on of order $(md)^3$ copies
of $G$; the hypothesis $(\mathrm M_{d,m}(c))$ is a bound on its second singular
value on gauge-invariant functions, uniform in the boundary condition, and the
argument above has no four-dimensional counterpart.

\subsection{Two dimensions: thin strips}\label{ssec:strips}

The restriction of $(\mathrm M_{d,m}(c))$ to boxes all of whose widths are at
least $2d-1$ is needed. We show that for strips of width one the same bound
fails along the two-dimensional tuning for every $c<1$.

Identify $\SU(2)$ with the unit quaternions $g=g_0+g_1\mathbf i+g_2\mathbf j+g_3
\mathbf k$; then $\chi:=\chi_{1/2}$ satisfies $\chi(g)=\Tr g=2g_0$, $\chi(g^{-1})
=\chi(g)$, and $\ip gh_{\R^4}=\frac12\chi(g^{-1}h)$. For $\epsilon\in\{1,-1\}$
and $a,b,c>0$ with $T=a+b+c$ let $\mu^\epsilon_{a,b,c}$ be the probability
measure on $\SU(2)^2$ with density
\[
p_a(x)\,p_b(x^{-1}y)\,p_c(\epsilon y)/p_T(\epsilon)
\]
with respect to $\lambda\otimes\lambda$: a Brownian bridge on $\SU(2)$ from $1$
to $\epsilon$ observed at the times $a$ and $a+b$. The normalisation follows
from $\int p_a(x)p_b(x^{-1}y)\,dx=p_{a+b}(y)$ and $\int p_{a+b}(y)p_c(\epsilon y)
\,dy=p_T(\epsilon)$, both instances of Lemma~\ref{lem:merge}(i).

\begin{lemma}[the strip is a bridge]\label{lem:striplaw}
Let $Q=[0,3d-1)\times\{0\}$, $d\ge2$, split in the first coordinate, let
$\varphi=p_t$, and let $\omega_\epsilon$ be the exterior configuration with all
links equal to $1$ except $\sigma_{((3d-1,0),2)}=\epsilon$. Then $\mathcal G_Q$
is trivial, so $\mathcal A^{\omega_\epsilon}_Q=L^2(\nu^{\omega_\epsilon}_Q)$,
and with $h_x=h(x,0)$ and $v_x=v(x,0)$ the pair $(x,y)=(v_{d-1}^{-1}h_{d-1},
v_{2d-1}^{-1})$, formed by a function of $\sigma_{V(Q_-)}$ and a function of
$\sigma_{V(Q_+)}$, has law $\mu^\epsilon_{a,b,c}$ with $a=c=2dt$ and
$b=(2d-1)t$.
\end{lemma}

\begin{proof}
No vertex $v$ of $Q$ has $v-e_2\in Q$, so there are no interior vertices. The
plaquettes meeting $V(Q)$ are those at $(-1,0)$, with holonomy $v_0$; at
$(x,-1)$, with holonomy $h_x^{-1}$; and at $(x,0)$, $0\le x\le3d-2$, with
holonomy $h_xv_{x+1}v_x^{-1}$, where $v_{3d-1}:=\epsilon$. So the density of
$\nu$ is proportional to $p_t(v_0)\prod_xp_t(h_x)\prod_xp_t(h_xv_{x+1}v_x^{-1})$.
For $x\ne d-1$, $\int p_t(h)p_t(hg)\,dh=(p_t*p_t)(g)=p_{2t}(g)$ by
Lemma~\ref{lem:merge}(i) and $p_t(h^{-1})=p_t(h)$. Integrating then
$v_0,\dots,v_{d-2}$, $v_{2d},\dots,v_{3d-2}$ and $v_d,\dots,v_{2d-2}$ with
Lemma~\ref{lem:merge}(i), in the form $\int\alpha(gy^{-1})\beta(yh)\,dy=(\alpha*
\beta)(gh)$, leaves the density of $(v_{d-1},h_{d-1},v_{2d-1})$ proportional to
\[
p_{(2d-1)t}(v_{d-1})\,p_t(h_{d-1})\,p_{(2d-1)t}(v_{2d-1}v_{d-1}^{-1}h_{d-1})\,
p_{2dt}(\epsilon v_{2d-1}^{-1}).
\]
Substituting $x=v_{d-1}^{-1}h_{d-1}$ for $h_{d-1}$, which preserves Haar measure
for fixed $v_{d-1}$, and integrating $v_{d-1}$ gives $(p_{(2d-1)t}*p_t)(x)=
p_{2dt}(x)$. With $y=v_{2d-1}^{-1}$ the third factor is $p_{(2d-1)t}(y^{-1}x)=
p_b(x^{-1}y)$, and the last is $p_{2dt}(\epsilon y)$.
\end{proof}

\begin{lemma}[moments of the bridge]\label{lem:moments}
Under $\mu^\epsilon_{a,b,c}$, $E[x_1]=E[y_1]=0$ and
\[
E[x_1y_1]=\tfrac13\bigl(\tfrac12D-\tfrac14C\bigr),\quad E[x_1^2]=\tfrac13
\bigl(1-\tfrac14A\bigr),\quad E[y_1^2]=\tfrac13\bigl(1-\tfrac14B\bigr),
\]
where $D=E\chi(x^{-1}y)$, $C=E[\chi(x)\chi(y)]$, $A=E\chi(x)^2$ and $B=E\chi(
y)^2$ are given, with $q_\theta(n)=e^{-\theta(n^2-1)/4}$ and $Z=p_T(\epsilon)$, by
\begin{align*}
Z&=\tfrac12\textstyle\sum_{n\in\Z}\epsilon^{n-1}n^2q_T(n),\qquad
DZ=\sum_{n\in\Z}\epsilon^{n-1}n(n+1)q_{a+c}(n)q_b(n+1),\\
CZ&=\textstyle\sum_{n\in\Z}\epsilon^n\bigl[(n+1)^2q_a(n+1)q_c(n+1)+(n^2-1)
q_a(n+1)q_c(n-1)\bigr]q_b(n),\\
AZ&=2Z+\textstyle\sum_{n\in\Z}\epsilon^{n-1}n(n+2)q_a(n+2)q_{b+c}(n),\\
BZ&=2Z+\textstyle\sum_{n\in\Z}\epsilon^{n-1}n(n+2)q_c(n+2)q_{a+b}(n).
\end{align*}
\end{lemma}

\begin{proof}
\emph{Symmetry.} Conjugation $(x,y)\mapsto(gxg^{-1},gyg^{-1})$ preserves
$\mu^\epsilon$, since $\epsilon$ is central and the $p_\theta$ are class
functions. It fixes $x_0$ and $y_0$ and rotates $(x_1,x_2,x_3)$ and
$(y_1,y_2,y_3)$ by the same element of $\SO(3)$. Conjugation by $\mathbf j$
maps $\mathbf i$ to $-\mathbf i$, so $E[x_1]=E[y_1]=0$; conjugation by
$\frac12(1+\mathbf i+\mathbf j+\mathbf k)$ permutes $\mathbf i,\mathbf j,
\mathbf k$ cyclically, so $E[x_1y_1]=\frac13E\sum_kx_ky_k$ and $E[x_1^2]=\frac13
E\sum_kx_k^2$. Finally $\sum_kx_ky_k=\ip xy-x_0y_0=\frac12\chi(x^{-1}y)-\frac14
\chi(x)\chi(y)$ and $\sum_kx_k^2=1-\frac14\chi(x)^2$.

\emph{Reduction to character sums.} Write class functions as $\sum_{n\ge1}F_n
\chi_{(n-1)/2}$. Then $p_\theta$ has coefficients $nq_\theta(n)$, and
$p_\theta(\epsilon\,\cdot)$ has coefficients $\epsilon^{n-1}nq_\theta(n)$, since
$\chi_j(\epsilon g)=\epsilon^{2j}\chi_j(g)$; moreover $\int(\sum F_n\chi)(\sum
G_n\chi)=\sum F_nG_n$, and $(\alpha*p_\theta*\gamma)(1)=\sum_n\alpha_n\gamma_n
q_\theta(n)$ because $\chi_i*\chi_j=\delta_{ij}\chi_j/d_j$ and $\chi_j(1)=d_j$.
By the Clebsch--Gordan rule $\chi\chi_j=\chi_{j-1/2}+\chi_{j+1/2}$ for $j\ge
\frac12$ and $\chi\chi_0=\chi$, multiplication by $\chi$ acts on coefficients
by $F_n\mapsto F_{n+1}+F_{n-1}$ with $F_0:=0$. The marginal of $x$ is $p_a(x)
p_{b+c}(\epsilon x)/Z$ and that of $y$ is $p_{a+b}(y)p_c(\epsilon y)/Z$;
substituting $z=x^{-1}y$, the law of $z$ is $p_b(z)p_{a+c}(\epsilon z)/Z$; and
$C=((\chi p_a)*p_b*\gamma)(1)/Z$ with $\gamma(w)=\chi(w)p_c(\epsilon w)$. All
series converge absolutely. Hence, with $P^{(\theta)}_n=nq_\theta(n)$ and
$\alpha^{(\theta)}_n=P^{(\theta)}_{n+1}+P^{(\theta)}_{n-1}$,
\begin{gather*}
Z=\sum_{n\ge1}\epsilon^{n-1}n^2q_T(n),\quad DZ=\sum_{n\ge1}\epsilon^{n-1}
\alpha^{(b)}_n\,nq_{a+c}(n),\quad CZ=\sum_{n\ge1}\epsilon^n\alpha^{(a)}_n
\alpha^{(c)}_nq_b(n),\\
AZ=\sum_{n\ge1}\epsilon^{n-1}(\alpha^{(a)}_{n+1}+\alpha^{(a)}_{n-1})\,nq_{b+c}
(n),\quad BZ=\sum_{n\ge1}\epsilon^{n-1}(\alpha^{(c)}_{n+1}+\alpha^{(c)}_{n-1})\,
nq_{a+b}(n),
\end{gather*}
where for $CZ$ we used that the coefficients of $\gamma$ are $\sum_\pm
\epsilon^{n\pm1-1}P^{(c)}_{n\pm1}=\epsilon^n\alpha^{(c)}_n$, and $\alpha_0:=0$.

\emph{Symmetrisation.} Extend $P^{(\theta)}_n=nq_\theta(n)$ to $n\in\Z$; it is
odd, so $\alpha^{(\theta)}$, now defined for all $n$, is odd with
$\alpha^{(\theta)}_0=0$ and agrees with the above for $n\ge1$, and so is
$\alpha_{n+1}+\alpha_{n-1}$. Each summand above is then an even function of $n$
vanishing at $0$, as $\epsilon^{n\pm1}$ is even in $n$, so $\sum_{n\ge1}=\frac12
\sum_{n\in\Z}$. Expanding the products and using the substitution $n\mapsto-n$,
under which $P_{n-1}\mapsto-P_{n+1}$ and $P_{n-2}\mapsto-P_{n+2}$, to identify
the terms in pairs gives the stated formulas; for $AZ$ the middle term
$2P^{(a)}_n\cdot nq_{b+c}(n)$ gives $2Z$, since $q_aq_{b+c}=q_T$.
\end{proof}

\begin{lemma}[Poisson summation]\label{lem:poisson}
Let $0<T\le1$, $\lambda=T/4$, $\gamma=\sqrt{4\pi/T}$, $\mu\in\R$, $\widetilde
Q(m)=q_2m^2+q_1m+q_0$ real, and $f(x)=\widetilde Q(x+\mu)e^{-\lambda(x+\mu)^2}$.
\begin{enumerate}[label=\textup{(\roman*)},leftmargin=*]
\item $|\sum_{n\in\Z}f(n)-\gamma(q_0+2q_2/T)|\le30\gamma(|q_0|+4|q_1|/T+16|q_2|/
T^2)e^{-4\pi^2/T}$.
\item If $q_2=1$, $|q_1|\le2$ and $|q_0|\le1$, then $\sum_{n\in\Z}(-1)^{n-1}f(n)
=\Lambda_0(\cos\pi\mu-\mathcal E)$ with $\Lambda_0=\frac{8\pi^2}{T^2}\gamma
e^{-\pi^2/T}$ and $|\mathcal E|\le0.22\,T$.
\end{enumerate}
\end{lemma}

\begin{proof}
$f$ is a Schwartz function. With $\widehat f(\xi)=\int f(x)e^{-2\pi i\xi x}dx$,
the substitution $y=x+\mu$ and the formulas $\int e^{-\lambda y^2-i\omega y}dy=
\gamma e^{-\omega^2/(4\lambda)}=:G(\omega)$, $\int ye^{-\lambda y^2-i\omega y}dy
=iG'(\omega)$, $\int y^2e^{-\lambda y^2-i\omega y}dy=-G''(\omega)$, $G'=-\frac
\omega{2\lambda}G$ and $G''=(\frac{\omega^2}{4\lambda^2}-\frac1{2\lambda})G$ give
\[
\widehat f(\xi)=e^{2\pi i\xi\mu}\gamma e^{-\pi^2\xi^2/\lambda}\Bigl[q_0-iq_1
\frac{\pi\xi}\lambda+q_2\Bigl(\frac1{2\lambda}-\frac{\pi^2\xi^2}{\lambda^2}
\Bigr)\Bigr].
\]
(i) By the Poisson summation formula~\cite[Ch.~5]{SS03}, $\sum_nf(n)=\sum_k
\widehat f(k)$, and $\widehat f(0)=\gamma(q_0+2q_2/T)$. For $k\ne0$, since
$\lambda\le\frac14$, each of $1$, $\pi|k|/\lambda$ and $\frac1{2\lambda}+
\frac{\pi^2k^2}{\lambda^2}$ is at most $M_k\lambda^{-e}$ with $e=0,1,2$
respectively and $M_k=1+\pi|k|+\pi^2k^2$; so $|\widehat f(k)|\le\gamma e^{-\pi^2
k^2/\lambda}M_k(|q_0|+|q_1|/\lambda+|q_2|/\lambda^2)$, and $\sum_{k\ne0}e^{-\pi^2
k^2/\lambda}M_k\le2e^{-\pi^2/\lambda}(M_1+\sum_{k\ge2}M_ke^{-12\pi^2(k-1)})\le
30e^{-\pi^2/\lambda}$, using $k^2-1\ge3(k-1)$ and $M_1<14.1$. Put $\lambda=T/4$.

(ii) Poisson summation for $f(x)e^{i\pi x}$ gives $\sum_n(-1)^nf(n)=\sum_{\xi\in
\frac12+\Z}\widehat f(\xi)$. The terms $\xi=\pm\frac12$ are complex conjugate
and add up to $2\gamma e^{-\pi^2/(4\lambda)}[\cos(\pi\mu)(q_0+\frac1{2\lambda}-
\frac{\pi^2}{4\lambda^2})+\sin(\pi\mu)\frac{q_1\pi}{2\lambda}]$. Hence
$\sum(-1)^{n-1}f(n)=\Lambda_0(\cos\pi\mu-\mathcal E)$ with $\Lambda_0=\frac{\pi^2}
{2\lambda^2}\gamma e^{-\pi^2/(4\lambda)}$, the stated value, and $\mathcal E=
\mathcal E_1+\mathcal E_2$, where $\mathcal E_1=\frac{2\lambda^2}{\pi^2}[\cos(\pi
\mu)(q_0+\frac1{2\lambda})+\sin(\pi\mu)\frac{q_1\pi}{2\lambda}]$ and $\mathcal
E_2=\Lambda_0^{-1}\sum_{|\xi|\ge3/2}\widehat f(\xi)$. First, $|\mathcal E_1|\le
\frac2{\pi^2}(\lambda^2+\frac\lambda2+\pi\lambda)\le\frac2{\pi^2}(\frac14+\frac12+
\pi)\lambda<0.79\lambda$. Second, for $|\xi|\ge\frac32$, $|\widehat f(\xi)|/
\Lambda_0\le e^{-\pi^2(\xi^2-1/4)/\lambda}(\frac{2\lambda^2}{\pi^2}+\frac{4\lambda
|\xi|}\pi+\frac\lambda{\pi^2}+2\xi^2)\le3\xi^2e^{-\pi^2(\xi^2-1/4)/\lambda}$, and
with $\xi=\pm(k+\frac12)$, $\xi^2-\frac14=k^2+k\ge2k$, so $|\mathcal E_2|\le6
\sum_{k\ge1}(k+\frac12)^2e^{-2\pi^2k/\lambda}\le14e^{-2\pi^2/\lambda}\le56
e^{-8\pi^2}\lambda$, as $e^{-2\pi^2/\lambda}/\lambda$ increases on $(0,\frac14]$.
So $|\mathcal E|\le0.8\lambda\le0.22\,T$.
\end{proof}

\begin{lemma}[the bridge from $1$ to $1$]\label{lem:bridge1}
Let $0<T\le1$ and suppose that $\kappa:=ac/T$, $\tau_x:=a(b+c)/T$ and $\tau_y:=
c(a+b)/T$ are all at least $T/10$. Under $\mu^1_{a,b,c}$,
\[
\operatorname{Corr}(x_1,y_1)\ge r_G(1-T),\qquad r_G=\Bigl(\frac{ac}{(a+b)(b+c)}
\Bigr)^{1/2}.
\]
\end{lemma}

Here $r_G$ is the correlation of a Euclidean Brownian bridge on $[0,T]$ at the
times $a$ and $a+b$.

\begin{proof}
\emph{Main terms.} Complete the squares in the exponents of
Lemma~\ref{lem:moments} with $\epsilon=1$: $q_T(n)=e^{T/4}e^{-Tn^2/4}$;
$q_{a+c}(n)q_b(n+1)=e^{E_D}e^{-T(n+\mu_D)^2/4}$ with $\mu_D=b/T$ and $E_D=
\frac{b^2}{4T}+\frac{a+c}4$; $q_a(n+1)q_c(n+1)q_b(n)=e^{E_1}e^{-T(n+\mu_1)^2/4}$
with $\mu_1=\frac{a+c}T$ and $E_1=\frac{(a+c)^2}{4T}+\frac b4$; $q_a(n+1)q_c(n-1)
q_b(n)=e^{E_2}e^{-T(n+\mu_2)^2/4}$ with $\mu_2=\frac{a-c}T$ and $E_2=\frac{(a-c)^2}
{4T}+\frac b4$; and $q_a(n+2)q_{b+c}(n)=e^{E_A}e^{-T(n+\mu_A)^2/4}$ with $\mu_A=
\frac{2a}T$ and $E_A=\frac{a^2}T-\frac{3a-b-c}4$. In the variable $m=n+\mu$ the
polynomial factors become $m^2$ for $Z$, $m^2+(1-2\mu_D)m+\mu_D^2-\mu_D$ for
$D$, $(m+1-\mu_1)^2$ with $1-\mu_1=b/T$ and $(m-\mu_2)^2-1$ for $C$, and
$(m-\mu_A)(m-\mu_A+2)$ for $A$. The main terms $\gamma(q_0+2q_2/T)$ of
Lemma~\ref{lem:poisson}(i), divided by the main term $\gamma e^{T/4}/T$ of $Z$,
give, after the simplifications $T-b=a+c$, $(a-c)^2-T(a+c)=-4ac-b(a+c)$ and
$a^2/T-a=-\tau_x$,
\begin{gather*}
D_0=e^{-u}(2-4u),\qquad C_0=e^{-u}\Bigl(2+\frac{b^2}T\Bigr)+e^{-u-\kappa}
\Bigl(2+\frac{(a-c)^2}T-T\Bigr),\\
A_0=2+e^{-\tau_x}(2-4\tau_x),\qquad B_0=2+e^{-\tau_y}(2-4\tau_y),\qquad
u=\frac{b(a+c)}{4T}.
\end{gather*}
Put $P=\frac{b(T+a+c)}{4T}$. Using $\frac{(a-c)^2-T^2}{4T}=-\frac{(2a+b)(2c+b)}
{4T}$ and $\frac{(2a+b)(2c+b)}{4T}=P+\kappa$, one finds
\[
X_0:=\tfrac12D_0-\tfrac14C_0=e^{-u}\bigl[(\tfrac12-P)(1-e^{-\kappa})+\kappa
e^{-\kappa}\bigr],\qquad V^x_0:=1-\tfrac14A_0=\tfrac12(1-e^{-\tau_x})+\tau_x
e^{-\tau_x}.
\]
Since $ac\le T^2/4$ and $b(a+c)\le T^2/4$, we have $\kappa\le T/4$, $u\le T/16$
and $P\le b/2\le T/2\le\frac12$. With $1-e^{-\kappa}\ge\kappa-\kappa^2/2$,
$e^{-\kappa}\ge1-\kappa$ and $e^{-u}\ge1-u$,
\[
X_0\ge(1-u)\kappa\bigl(\tfrac32-P-\tfrac54\kappa\bigr)\ge\tfrac32\kappa\bigl(1-
\tfrac T{16}\bigr)\bigl(1-\tfrac{13}{24}T\bigr)\ge\tfrac32\kappa\bigl(1-\tfrac{29}
{48}T\bigr).
\]
Also $V^x_0\le\frac32\tau_x$ and $V^y_0:=1-\frac14B_0\le\frac32\tau_y$.

\emph{Errors.} Each of the five sums is its main term times $1+\theta$ with
$|\theta|\le\eta:=750e^{-4\pi^2/T}/T$. This follows from
Lemma~\ref{lem:poisson}(i), with $T\le1$, the coefficient bounds $q_2=1$,
$|q_1|\le2$ and $|q_0|\le1$ read off above, and the lower bounds $2/T-1/4$,
$2/T$, $2/T-1$ and $2/T-1\ge1/T$ for the brackets $q_0+2q_2/T$ of the main
terms: the relative errors are at most $240e^{-4\pi^2/T}/T$ for $Z$ and $750
e^{-4\pi^2/T}/T$ for the others. As $e^{-4\pi^2/T}/T^3$ increases on $(0,1]$,
$\eta\le750e^{-4\pi^2}T^2<10^{-14}T^2$. Since $|D_0|\le2$, both parts of $C_0$
are at most $3$ in absolute value, and $0\le A_0-2\le2$, we get $|D-D_0|\le6
\eta$, $|C-C_0|\le18\eta$ and $|A-A_0|,|B-B_0|\le6\eta$, hence $|X-X_0|\le7.5
\eta$ and $|V^x-V^x_0|,|V^y-V^y_0|\le1.5\eta$, where $X=\frac12D-\frac14C$,
$V^x=1-\frac14A$ and $V^y=1-\frac14B$.

\emph{Conclusion.} By Lemma~\ref{lem:moments}, $\operatorname{Corr}(x_1,y_1)=X/
(V^xV^y)^{1/2}$. Using $\kappa,\tau_x,\tau_y\ge T/10$,
\[
\operatorname{Corr}(x_1,y_1)\ge\frac{\frac32\kappa(1-\frac{29}{48}T-50\eta/T)}
{\frac32(\tau_x\tau_y)^{1/2}(1+10\eta/T)}\ge r_G\Bigl(1-\frac{29}{48}T-60\frac
\eta T\Bigr)\ge r_G(1-T),
\]
since $\kappa/(\tau_x\tau_y)^{1/2}=r_G$ and $60\eta/T<10^{-12}T$.
\end{proof}

\begin{lemma}[the bridge from $1$ to $-1$]\label{lem:bridgem1}
Let $s_0>0$ and $0<T\le\min(1,s_0^2/4)$, and suppose $\sin(\pi a/T)\ge s_0$ and
$\sin(\pi c/T)\ge s_0$. Under $\mu^{-1}_{a,b,c}$,
\[
\operatorname{Corr}(x_1,y_1)\ge1-\frac{5T}{2s_0^2}.
\]
\end{lemma}

\begin{proof}
Write $r_a=a/T$, $r_b=b/T$, $r_c=c/T$, $s_a=\sin\pi r_a$ and $s_c=\sin\pi r_c$.
The sums of Lemma~\ref{lem:moments} with $\epsilon=-1$ have the exponents and
polynomial factors of the proof of Lemma~\ref{lem:bridge1}, all with $q_2=1$,
$|q_1|\le2$ and $|q_0|\le1$, and carry the sign $(-1)^{n-1}$, except $CZ$, which
carries $(-1)^n$. Lemma~\ref{lem:poisson}(ii) gives
\begin{gather*}
Z=\tfrac12e^{T/4}\Lambda_0(1-\mathcal E_Z),\qquad D=2e^{-u}\frac{\cos\pi r_b-
\mathcal E_D}{1-\mathcal E_Z},\\
C=-2\,\frac{e^{-u}(\cos\pi(r_a+r_c)-\mathcal E_1)+e^{-u-\kappa}(\cos\pi(r_a-r_c)
-\mathcal E_2)}{1-\mathcal E_Z},\\
A=2+2e^{-\tau_x}\frac{\cos2\pi r_a-\mathcal E_A}{1-\mathcal E_Z},\qquad
B=2+2e^{-\tau_y}\frac{\cos2\pi r_c-\mathcal E_B}{1-\mathcal E_Z},
\end{gather*}
with all $|\mathcal E_\bullet|\le0.22T$ and $u,\kappa,\tau_x,\tau_y$ as before;
the exponential prefactors are those of the main terms in the proof of
Lemma~\ref{lem:bridge1}. Each quotient has the form $w(\cos\theta-\mathcal E)/(1-
\mathcal E_Z)$ with $w\in[1-T/4,1]$, as $u+\kappa,\tau_x,\tau_y\le T/4$; since
$|1/(1-\mathcal E_Z)|\le1.29$ and $|1/(1-\mathcal E_Z)-1|\le0.29T$, it differs
from $\cos\theta$ by at most $0.25T\cdot1.29+0.29T+0.22T\cdot1.29<T$. Hence
\[
\bigl|\tfrac D2-\cos\pi r_b\bigr|\le T,\quad\bigl|\tfrac C4+\cos\pi r_a\cos\pi
r_c\bigr|\le T,\quad\bigl|\tfrac A4-\cos^2\pi r_a\bigr|\le\tfrac T2,\quad\bigl|
\tfrac B4-\cos^2\pi r_c\bigr|\le\tfrac T2,
\]
using $\frac12(\cos\pi(r_a+r_c)+\cos\pi(r_a-r_c))=\cos\pi r_a\cos\pi r_c$. Since
$r_b=1-r_a-r_c$, $\cos\pi r_b=-\cos\pi(r_a+r_c)$, so $\cos\pi r_b+\cos\pi r_a
\cos\pi r_c=s_as_c$. Thus $X=\frac12D-\frac14C\ge s_as_c-2T>0$, and $0<V^x=1-
\frac14A\le s_a^2+\frac T2$ and $0<V^y\le s_c^2+\frac T2$, positivity coming
from Lemma~\ref{lem:moments} and $T\le s_0^2/4$. By Lemma~\ref{lem:moments} and
$\sqrt{(1+x)(1+y)}\le1+\frac{x+y}2$,
\[
\operatorname{Corr}(x_1,y_1)\ge\frac{s_as_c-2T}{s_as_c(1+\frac T{4s_a^2}+\frac T
{4s_c^2})}\ge\frac{1-2T/s_0^2}{1+T/(2s_0^2)}\ge1-\frac{5T}{2s_0^2}.
\qedhere
\]
\end{proof}

\begin{proposition}[thin strips]\label{prop:strip}
Let $G=\SU(2)$, $\varphi=p_t$, $d\ge2$, let $Q=[0,3d-1)\times\{0\}$ be split in
the first coordinate into widths $d$, $d-1$, $d$, put $T=(6d-1)t$, and let
$\omega_{\pm1}$ be as in Lemma~\ref{lem:striplaw}.
\begin{enumerate}[label=\textup{(\roman*)},leftmargin=*]
\item If $T\le1$, then $\rho^{\omega_1}(Q_-,Q_+)\ge\frac{2d}{4d-1}(1-T)$.
\item If $T\le\frac3{16}$, then $\rho^{\omega_{-1}}(Q_-,Q_+)\ge1-\frac{10}3T$.
\item Along the tuning $t_k=g^2a_k^2$, $a_k\to0$, with integers $d_k\ge2$ and
$d_ka_k\le\ell$, for every $m\ge4$ and every $c<1$ the bound $\rho^\omega(Q_-,
Q_+)\le c$ fails for all large $k$ for the strip of widths $3d_k-1$ and $1$,
both at most $md_k-1$; already the trivial boundary condition gives $\rho\ge
\frac12(1-6g^2\ell a_k)$.
\end{enumerate}
\end{proposition}

\begin{proof}
By Lemma~\ref{lem:striplaw} the functions $X=x_1$ and $Y=y_1$ of
$\sigma_{V(Q_-)}$ and $\sigma_{V(Q_+)}$ lie in $\mathcal A^{\omega_\epsilon}_Q$,
and their joint law is that of $\mu^\epsilon_{a,b,c}$ with $a=c=2dt$ and $b=(2d-
1)t$; so $\rho^{\omega_\epsilon}\ge|\operatorname{Corr}(x_1,y_1)|$. (i) Here
$\kappa/T=4d^2/(6d-1)^2\ge\frac19$, $\tau_x/T=\tau_y/T=2d(4d-1)/(6d-1)^2\ge\frac15$
for $d\ge2$, and $r_G=2d/(4d-1)$; apply Lemma~\ref{lem:bridge1}. (ii) $a/T=c/T=
2d/(6d-1)\in(\frac13,\frac12)$, so $\sin(\pi a/T)\ge\sin\frac\pi3=\frac{\sqrt3}2=
:s_0$ with $s_0^2/4=\frac3{16}$; apply Lemma~\ref{lem:bridgem1}. (iii)
$T_k=(6d_k-1)t_k\le6g^2(d_ka_k)a_k\le6g^2\ell a_k\to0$, so for large $k$, by
(ii), $\rho^{\omega_{-1}}\ge1-\frac{10}3T_k>c$, and by (i), $\rho^{\omega_1}\ge
\frac{2d_k}{4d_k-1}(1-T_k)\ge\frac12(1-6g^2\ell a_k)$.
\end{proof}

The mechanism is a bridge. The fixed boundary pins the holonomy of the loop
surrounding the strip, and both outer parts read off the same collective
variable. For the condition $\omega_{-1}$ the measure is invariant under the
global conjugation $\sigma\mapsto g\sigma g^{-1}$ of all links of $V(Q)$, the
exterior links being central, and part (ii) shows that this global $\SU(2)$
symmetry is broken at small $T$: the bridge from $1$ to $-1$ concentrates near
the two-sphere of minimising geodesics, whose point is read off by both outer
parts, and the slab correlation tends to one. Theorem~\ref{thm:twod} and
Proposition~\ref{prop:strip} together show that in two dimensions, along the
weak-coupling tuning of a non-abelian theory, the slab hypothesis holds for fat
boxes uniformly in the boundary condition with $c\to0$, and fails for strips
of width one for every $c<1$. In four dimensions the boundary conditions with
values in the centre are 't~Hooft's twisted ones~\cite{tH79}, and for fat boxes
they are the natural test of $(\mathrm M_{d,m}(c))$; we have no estimate for
them along the four-dimensional tuning.

\subsection{Curvature}\label{ssec:curvlimits}

Theorem~\ref{thm:A} is a curvature argument on the gauge-fixed measure, and its
range is limited by the Hessian bound of Theorem~\ref{thm:hess}. We show that
no curvature argument applied to the Wilson measure itself, with or without the
gauge, can be uniform along the weak-coupling tuning. We use the notation of
Section~\ref{sec:curvature}: $G=\SU(N)$ with the metric $\ip XY=-\Tr(XY)$,
tangent vectors $v=(X_eQ_e)$, the transported vectors $Y_1,\dots,Y_4$ of a
plaquette, and $S_0=\sum_p\ReTr Q_p$. For a probability measure with density
proportional to $e^{-V}$ on a compact Riemannian manifold, the
\emph{Bakry--\'Emery constant} is $\inf(\Ric+\Hess V)(v,v)/|v|^2$ over all
points and nonzero tangent vectors; by Bochner's formula it is the largest $K$
with $\Gamma_2(f)\ge K\Gamma(f)$ pointwise for all smooth $f$, and it controls
the Poincar\'e and logarithmic Sobolev constants of the Langevin
dynamics~\cite{BGL14}. For a plaquette put $Z_p=Y_1+Y_2-Y_3-Y_4$.

\begin{lemma}[central plaquettes]\label{lem:central}
Let $Q$ be a field and $p$ a plaquette with $Q_p=wI$ central. Along the
geodesics $Q_e(t)=e^{tX_e}Q_e$,
\[
\frac{d^2}{dt^2}\Big|_{t=0}\ReTr Q_p(t)=-\operatorname{Re}(w)\,|Z_p|^2 .
\]
If $\varphi$ is a smooth class function on $G$, then $\frac{d^2}{dt^2}|_0
\varphi(Q_p(t))=H_w(Z_p,Z_p)$ for an $\operatorname{Ad}$-invariant quadratic
form $H_w$ on $\g$; for $G=\SU(2)$, $H_{\pm1}=h_\pm|\cdot|^2$ with real constants
$h_\pm$.
\end{lemma}

\begin{proof}
As in Section~\ref{sec:curvature}, $Q_p(t)=e^{tY_1}e^{tY_2}We^{-tY_3}e^{-tY_4}$
with $W=Q_p$. Since $W=wI$ is central, $Q_p(t)=w\,c(t)$ with $c(t)=e^{tY_1}
e^{tY_2}e^{-tY_3}e^{-tY_4}$. With $a=(1,1,-1,-1)$, $c''(0)=\sum_ia_i^2Y_i^2+2
\sum_{i<j}a_ia_jY_iY_j$, and as $\Tr(Y_iY_j)=\Tr(Y_jY_i)$, $\Tr c''(0)=\Tr(Z_p^2)
=-|Z_p|^2$. Hence $\frac{d^2}{dt^2}\ReTr Q_p=\operatorname{Re}(w\Tr Z_p^2)=
-\operatorname{Re}(w)|Z_p|^2$, the trace of a product of two skew-Hermitian
matrices being real. For $\varphi$, the curve $\gamma(t)=Q_p(t)$ has
$\gamma(0)=wI$ and $\gamma'(0)=Z_pwI$, and $(\varphi\circ\gamma)''(0)=\Hess
\varphi(\gamma'(0),\gamma'(0))+d\varphi(\nabla_t\gamma'(0))$. The gradient of
$\varphi$ at $wI$ is $YwI$ with $Y$ fixed by every $\operatorname{Ad}_g$, because
conjugation by $g$ is an isometry that fixes $wI$ and preserves $\varphi$; the
centre of $\mathfrak{su}(N)$ is trivial, so $Y=0$ and the second term vanishes.
The same isometries show that $H_w(Z,Z)=\Hess\varphi_{wI}(ZwI,ZwI)$ is
$\operatorname{Ad}$-invariant, and $\operatorname{Ad}$ is irreducible on
$\mathfrak{su}(2)$, so $H_{\pm1}$ is a multiple of the metric.
\end{proof}

The configuration
\begin{equation}\label{eq:piflux}
U^\pi_{x,\mu}=(-1)^{\sum_{\nu<\mu}x_\nu}I
\end{equation}
has every plaquette equal to $-I$: for $\mu<\nu$ the link $(x+e_\mu,\nu)$ sees
the coordinate $x_\mu$ raised by one and the link $(x+e_\nu,\mu)$ does not, so
the four signs of the plaquette at $x$ multiply to $-1$. It is defined on every
box and on $\Lambda_L$ for even $L$, and lies in $\SU(N)$ for even $N$. Put
$s(x)=(-1)^{|x|}$ with $|x|=\sum_\mu x_\mu$.

\begin{theorem}[full space]\label{thm:fullspace}
Let $N$ be even, $L\ge4$ even, and $\mu_L$ the $\SU(N)$ Wilson measure on
$\Lambda_L$ with density proportional to $e^{N\beta S_0}$. Its Bakry--\'Emery
constant satisfies
\[
K_{\rm full}\le\frac N2-16N|\beta|=\frac N2-\frac{16|\beta_W|}N .
\]
For $\SU(2)$, $K_{\rm full}\le1-8|\beta_W|$.
\end{theorem}

\begin{proof}
By Lemma~\ref{lem:ricci} and Table~\ref{tab:groups}, $\Ric=\frac N2\ip\cdot\cdot$.
Let $c\in\R^4$ with $\sum_\mu c_\mu=0$ and $c\ne0$, let $E\in\mathfrak{su}(N)$
with $|E|=1$, and put $X_{x,\mu}=c_\mu s(x)E$. For $\beta\ge0$ take the field
$U^\pi$, at which all transports are $\operatorname{Ad}_{\pm I}=\mathrm{id}$.
For the plaquette at $x$ in the plane $\mu<\nu$, $Y_1=c_\mu s(x)E$,
$Y_2=-c_\nu s(x)E$, $Y_3=-c_\mu s(x)E$ and $Y_4=c_\nu s(x)E$, so $Z_p=2s(x)(c_\mu-
c_\nu)E$. By Lemma~\ref{lem:central} with $w=-1$,
\[
\Hess S_0(v,v)=\sum_p|Z_p|^2=4L^4\sum_{\mu<\nu}(c_\mu-c_\nu)^2=4L^4\Bigl(4\sum_\mu
c_\mu^2-\Bigl(\sum_\mu c_\mu\Bigr)^2\Bigr)=16|v|^2 .
\]
With $V=-N\beta S_0$ this gives $(\Ric+\Hess V)(v,v)=(\frac N2-16N\beta)|v|^2$.
For $\beta<0$ take the field $Q\equiv I$, where Lemma~\ref{lem:central} with
$w=1$ gives $\Hess S_0(v,v)=-16|v|^2$ for the same $v$, hence $(\Ric+\Hess V)(v,
v)=(\frac N2-16N|\beta|)|v|^2$. Finally $\beta_W=N^2\beta$, and $\beta_W=4\beta$
for $\SU(2)$.
\end{proof}

Without the gauge the bound $|\Hess S_0(v,v)|\le16|v|^2$ holds for every field
and vector~\cite{BB26a}, so Theorem~\ref{thm:fullspace} is an equality,
$K_{\rm full}=\frac N2-16N|\beta|$. The Ricci curvature grows with $N$, but so
does the Hessian, and the criterion $K_{\rm full}>0$ holds exactly for
$|\beta_W|<N^2/32$; at fixed $N$ along the tuning $\beta_W\to\infty$ the
constant tends to $-\infty$ linearly. The Bakry--\'Emery constant is also the
best rate of contraction of the Langevin dynamics in the quadratic Kantorovich
distance~\cite{vRS05}, so no synchronous coupling of the dynamics does better.
The gauge of Section~\ref{sec:linefield} improves the constant $16$ to
$\frac{29}2$, and the next theorem shows that it cannot improve it below
$\frac{38}3$.

\begin{theorem}[the gauge slice]\label{thm:slice}
Let $G=\SU(2)$, $m\ge1$, $\Lambda_m=[0,2m]^4\cap\Z^4$ with free boundary, and
let $K_F(\Lambda_m)$ be the Bakry--\'Emery constant of the gauge-fixed measure
$\mu^F_{\Lambda_m}$ of Lemma~\ref{lem:gaugefix}. Then
\[
K_F(\Lambda_m)\le1-2|\beta|\rho_m,\qquad\rho_m=\frac{304\,m^4}{24\,m^4+8(2m+1)^3}
\xrightarrow[m\to\infty]{}\frac{38}3 .
\]
Hence the constant $\frac{29}2$ of Theorem~\ref{thm:hess} cannot be replaced by
any number below $\frac{38}3$, and asymptotically in $m$
\[
1-\tfrac{29}4|\beta_W|\le K_F\le1-\tfrac{19}3|\beta_W| ,
\]
so no Bakry--\'Emery argument for the gauge-fixed measure reaches $|\beta_W|=
\frac3{19}$.
\end{theorem}

\begin{proof}
\emph{The field.} For $\beta\ge0$ start from $U^\pi$ on $\Lambda_m$ and apply
the gauge transformation of the proof of Lemma~\ref{lem:gaugefix} along the
forest $F\cap E(\Lambda_m)$; the recursion $g_{x_{i+1}}=g_{x_i}U_{(x_i,x_{i+1})}$
produces $g_x\in\{\pm I\}$, so the result $Q$ takes values in $\{\pm I\}$,
equals $I$ on $F$, and has all plaquettes conjugate to those of $U^\pi$, that is
$-I$. For $\beta<0$ take $Q\equiv I$. For $v$ vanishing on $F$ the Hessian of
the gauge-fixed action is $\Hess S_{\Lambda_m}(v,v)$, which by
Lemma~\ref{lem:central} equals $\pm\sum_{p\in P(\Lambda_m)}|Z_p|^2$, with $+$ for
$\beta\ge0$ and $-$ for $\beta<0$; all transports are trivial.

\emph{The vector.} Let $c=(1,-1,0,0)$, $|E|=1$, and put $X_{x,\mu}=c_\mu s(x)E$
on the active links of $E(\Lambda_m)$ and $X=0$ on $F$. Let $p$ be a plaquette at
$x$ in the plane $\mu<\nu$ with all its links in $E(\Lambda_m)$; it contains
exactly one link of $F$. Without that link the computation of the proof of
Theorem~\ref{thm:fullspace} gives $Z_p=2s(x)(c_\mu-c_\nu)E$, the links $e_1,e_3$
of direction $\mu$ contributing $s(x)c_\mu E$ each and the links $e_2,e_4$ of
direction $\nu$ contributing $-s(x)c_\nu E$ each. Removing the link of $F$
gives $Z_p=s(x)(c_\mu-2c_\nu)E$ if that link has direction $\mu$, and $Z_p=
s(x)(2c_\mu-c_\nu)E$ if it has direction $\nu$.

\emph{Counting.} Let $B=[0,2m-1]^4$; every plaquette with base point in $B$ has
its links in $E(\Lambda_m)$. Fix $\mu<\nu$. The link of $F$ of the plaquette at
$x$ has direction $\mu$ exactly when $\sigma(x)=\mu$ or $\sigma(x+e_\nu)=\mu$,
and the two cases exclude each other. By Table~\ref{tab:sigma}, $\sigma$ has
period two and takes each value at four of the sixteen residues. $B$ and
$B+e_\nu$ are disjoint unions of $m^4$ translates of $\{0,1\}^4$, each
containing every residue once, so each case occurs for $4m^4$ base points in
$B$; of the $16m^4$ plaquettes of the plane with base point in $B$, exactly
$8m^4$ have their link of $F$ in direction $\mu$ and $8m^4$ in direction $\nu$.
Hence, with $|c|^2=2$ and $\sum_\mu c_\mu=0$,
\[
\sum_{p\in P(\Lambda_m)}|Z_p|^2\ge8m^4\sum_{\mu<\nu}\bigl[(c_\mu-2c_\nu)^2+
(2c_\mu-c_\nu)^2\bigr]=8m^4\Bigl(19|c|^2-4\Bigl(\sum_\mu c_\mu\Bigr)^2\Bigr)=
304m^4,
\]
using $\sum_{\mu<\nu}(5c_\mu^2+5c_\nu^2-8c_\mu c_\nu)=15|c|^2-4((\sum_\mu c_\mu
)^2-|c|^2)$. On the other side, for $\mu\in\{0,1\}$ the links $(x,\mu)$ with
$x\in B$ lie in $E(\Lambda_m)$, and exactly $16m^4-4m^4=12m^4$ of them are
active; they contribute $24m^4$ to $|v|^2$. The remaining links of
$E(\Lambda_m)$ with $X\ne0$ have direction $0$ or $1$ and base point in
$[0,2m]^4\setminus B$, at most $(2m+1)^4-(2m)^4\le4(2m+1)^3$ in each direction,
and each contributes at most $1$. So $|v|^2\le24m^4+8(2m+1)^3$, and
$|\Hess S_{\Lambda_m}(v,v)|\ge\rho_m|v|^2$ with the sign that gives $\Hess V(v,
v)\le-2|\beta|\rho_m|v|^2$, $V=-2\beta S_{\Lambda_m}$.

\emph{Conclusion.} $(\Ric+\Hess V)(v,v)\le(1-2|\beta|\rho_m)|v|^2$. If
Theorem~\ref{thm:hess} held with a constant $c'<\frac{38}3$, then $\rho_m\le c'$
for all $m$, contradicting $\rho_m\to\frac{38}3$. The lower bound for $K_F$ is
Lemma~\ref{lem:weighted} with $\lambda=0$, and $\beta_W=4\beta$.
\end{proof}

\begin{proposition}[plaquette actions at a coarse scale]\label{prop:coarse}
Let $L\ge4$ be even and $\nu(dU)\propto\exp(\sum_p\phi(U_p))\prod_edU_e$ on
$\SU(2)^{E(\Lambda_L)}$ with $\phi(U)=g(\frac12\ReTr U)$, $g\in C^2([-1,1])$. The
Bakry--\'Emery constant $K$ of $\nu$ satisfies
\[
K\le1-8\max\bigl(g'(-1),-g'(1)\bigr).
\]
\end{proposition}

\begin{proof}
$\frac12\ReTr(\pm e^{tZ})=\pm(1-\frac{t^2}4|Z|^2)+O(t^3)$, since $\Tr Z^2=
-|Z|^2$; so in Lemma~\ref{lem:central} $h_+=-\frac12g'(1)$ and $h_-=\frac12
g'(-1)$. The potential is $-\sum_p\phi(U_p)$. At $U^\pi$, with the vector of
the proof of Theorem~\ref{thm:fullspace}, $\Hess\sum_p\phi(U_p)(v,v)=h_-\sum_p
|Z_p|^2=16h_-|v|^2$, and at $Q\equiv I$ it is $16h_+|v|^2$. Hence $K\le1-16
h_\pm$.
\end{proof}

For the Wilson weight $g(s)=4\beta s$ this is Theorem~\ref{thm:fullspace}. For
an effective plaquette action at a coarse scale the Bakry--\'Emery route needs
$|g'(\pm1)|<\frac18$, an effective theory deep in strong coupling, and producing
such an effective measure from the weak-coupling lattice theory, and
transporting the resulting estimate back, are the two missing steps of
Section~\ref{sec:weak}.

\subsection{Entropy, link covariances and blocks of fixed size}
\label{ssec:entropy}

In the setting of Section~\ref{sec:cube}, say that $\mu$ satisfies \emph{block
entropy factorisation} with constant $C$ if $\operatorname{Ent}_\mu(F)\le C
n^{-4}\sum_x\mu(\operatorname{Ent}_{\gamma^0_{V_x}}F)$ for all bounded
gauge-invariant $F\ge0$, where $\operatorname{Ent}_{\gamma^0_{V_x}}F$ is the
entropy of $F$ under $\gamma^0_{V_x}(\cdot\mid U)$.

\begin{proposition}\label{prop:entbp}
Block entropy factorisation with constant $C$ implies $\BP(n,1/C)$.
\end{proposition}

\begin{proof}
Let $\Phi(s)=s\log s$. For $|u|\le\frac12$, $\Phi(1+u)=u+\frac{u^2}2+r(u)$ with
$r(u)=\sum_{k\ge3}(-u)^k/(k(k-1))$ and $|r(u)|\le|u|^3\sum_{k\ge3}2^{3-k}/(k(k-1))
\le|u|^3$. Let $g$ be real, bounded and gauge invariant, and $0<\varepsilon\le1/
(2\|g\|_\infty)$. Then $\operatorname{Ent}_\mu(1+\varepsilon g)=\mu\Phi(1+
\varepsilon g)-\Phi(1+\varepsilon\mu g)=\frac{\varepsilon^2}2\Var_\mu g+R$ with
$|R|\le2\varepsilon^3\|g\|_\infty^3$, and in the same way
$\operatorname{Ent}_{\gamma^0_{V_x}}(1+\varepsilon g)=\frac{\varepsilon^2}2
\Var_{\gamma^0_{V_x}}g+R_x$ with $|R_x|\le2\varepsilon^3\|g\|_\infty^3$. Inserting
$F=1+\varepsilon g$, dividing by $\varepsilon^2/2$ and letting $\varepsilon\to0$
gives $\Var_\mu g\le Cn^{-4}\sum_x\mu(\Var_{\gamma^0_{V_x}}g)=Cn^{-4}\sum_x\mu((g-
\mathcal P_xg)^2)$. Truncations of $f\in\mathcal A$ are bounded and gauge
invariant and converge in $L^2(\mu)$, and both sides are continuous on
$L^2(\mu)$.
\end{proof}

So entropy factorisation at the physical scale, uniform along the tuning,
implies (D) with $\varpi=1/C$, and every obstruction to the cube condition
applies to it: it fails for $\U(1)$ at weak coupling and at fixed cube size
(Section~\ref{sec:weak}). The methods that derive block factorisation from
spatial mixing uniformly in the boundary condition~\cite{CP21} need, at the
physical scale, an input of the strength of (D) itself.

\begin{proposition}[link covariances]\label{prop:linkcov}
Let $\mu$ be a gauge-invariant probability measure on $\SU(2)^E$, $E$ the links
of a box or of $\Lambda_L$ with $L\ge3$, and identify $U_e$ with a unit vector
$\sigma_e\in\R^4$. Then $E\sigma_e=0$, $E\sigma_e\sigma_e^T=\frac14I_4$ and
$E\sigma_e\sigma_{e'}^T=0$ for $e\ne e'$.
\end{proposition}

\begin{proof}
Distinct links have distinct sets of end points, so $e$ has an end point $x$
that is not an end point of $e'$. The gauge transformation $g$ at $x$ alone
maps $\sigma_e$ to $\ell_g\sigma_e$, where $\ell_g$ is left multiplication by
$g$ or right multiplication by $g^{-1}$, an orthogonal map, and fixes
$\sigma_{e'}$. By invariance $E\sigma_e\sigma_{e'}^T=\int\ell_g\,dg\;E\sigma_e
\sigma_{e'}^T=0$, since $\ell_g$ is linear in $g$ and $\int g\,dg=0$. The same
argument gives $E\sigma_e=0$ and shows that the law of $\sigma_e$ is invariant
under multiplication, hence Haar measure, so $E\sigma_e\sigma_e^T=\frac14I_4$.
\end{proof}

The influence and covariance matrices that spectral independence and
localisation schemes ask to be bounded are therefore exactly $\frac14I$ for the
Wilson measure at every $\beta$, weak coupling included, and carry no
information. Those methods need the bound for all tilted measures, and the
tilts break gauge invariance. Their output is a tensorisation inequality over
blocks of bounded lattice size, uniform in the volume, and the next
proposition shows that such an inequality cannot hold uniformly along the
tuning when the continuum limit is nontrivial.

\begin{proposition}[fixed blocks force triviality]\label{prop:fixedblocks}
Let $a_k=a_02^{-k}$ and $\beta_k$ be given, and suppose that for every $k$
there is $C_k<\infty$ such that the Wilson measure $\mu_{\beta_k,L}$ satisfies
\[
\Var_\mu(f)\le C_k\sum_{e\in E}\mu(\Var_ef)\qquad(f\in\mathcal A)
\]
for all $L$ in an infinite set $\mathcal L_k$ of even integers $L\ge8$, where
$\Var_ef$ is the conditional variance given the links other than $e$. Then
$\BP(2,1/(64C_k))$ holds on these tori, every torus limit along $\mathcal L_k$
has transfer-matrix gap at least $1/(5000e^2C_k)$, and every continuum limit
of such torus limits in the sense of Definition~\ref{def:limit} has
$\sigma(H)\subset\{0\}\cup[E_*,\infty)$ with $E_*=\limsup_k(5000e^2C_ka_k)^{-1}$.
In particular, if $C_ka_k\to0$, every continuum limit is trivial.
\end{proposition}

\begin{proof}
For $e\in V_x$ the $\sigma$-algebra of the links off $V_x$ is contained in that
of the links other than $e$, so $\mu(\Var_ef)=\|f-E[f\mid e^c]\|^2\le\|f-
\mathcal P_xf\|^2$. With $n=2$ each link lies in $16$ cubes $V_x$ and $|V_x|=4
\cdot2^4=64$, so
\[
\sum_e\mu(\Var_ef)\le\frac1{16}\sum_x\sum_{e\in V_x}\mu((f-\mathcal P_xf)^2)=4
\sum_x\mu((f-\mathcal P_xf)^2),
\]
and $\Var_\mu f\le64C_k\cdot2^{-4}\sum_x\mu((f-\mathcal P_xf)^2)$, which is
$\BP(2,1/(64C_k))$ on tori with $L\ge8=2\cdot2+4$. Corollary~\ref{cor:latticegap}
gives the gap $m_k=\varpi/(2e^2q_2)$ with $q_2=\frac{625}{16}$, that is $m_k=1/
(5000e^2C_k)$. For a continuum limit, as in the proof of Theorem~\ref{thm:B},
for $f\in\mathcal T$, $\varphi=\psi(f)-\ip\Omega{\psi(f)}\Omega$ and dyadic
$\tau>0$, $\ip\varphi{e^{-\tau H}\varphi}=\lim_k\ip{\varphi_k}{\mathbb T_k^{\tau/
a_k}\varphi_k}\le\liminf_ke^{-m_k\tau/a_k}\|\varphi_k\|^2=e^{-E_*\tau}\|\varphi
\|^2$, using $\|\varphi_k\|^2\to\|\varphi\|^2$, the case $\tau=0$ of the
definition. These $\varphi$ are dense in $\Omega^\perp$, and the spectral
argument of the proof of Theorem~\ref{thm:B} gives $\sigma(H|_{\Omega^\perp})
\subset[E_*,\infty)$. If $E_*=\infty$, then $\ip\varphi{e^{-\tau H}\varphi}=0$
for all dyadic $\tau>0$, and strong continuity of $e^{-\tau H}$ at $\tau=0$
gives $\varphi=0$, so $\Omega^\perp=\{0\}$.
\end{proof}

The same holds for blocks of any fixed lattice size $b$ in place of single
links, with $\BP(b,\cdot)$. So the tensorisation constants of single links or of
bounded blocks must grow at least like $1/a_k$ along the tuning whenever the
continuum limit is nontrivial; a usable version must work with blocks of
physical size, and is then a form of (D).

\subsection{The transfer matrix alone}\label{ssec:transfer}

The next results show what the continuum gap needs and that a lattice gap
tending to zero at the right rate is, by itself, a kinematic statement. In this
section Definition~\ref{def:limit} is used with $a_k=a_02^{-j_k}$ for a strictly
increasing sequence of integers $j_k$, with $\tau$ in $\mathbb D=\{a_0p2^{-q}:
p,q\in\Z_{\ge0}\}$, and with an arbitrary complex vector space $\mathcal T$ and
linear maps $O_k$ into the gauge-invariant local observables at nonnegative
times; the proof of Theorem~\ref{thm:B} uses nothing more. The states
$\omega_k$ are translation invariant and reflection positive, with transfer
matrices $0\le\mathbb T_k\le1$ on $\mathcal H_k$, vacua $\Omega_k$ and maps
$\Psi_k$ from observables at nonnegative times onto dense subspaces, as in
Section~\ref{sec:lattice}. Put $\gamma_k=\|\mathbb T_k|_{\Omega_k^\perp}\|$,
$m_k=-\log\gamma_k\in[0,\infty]$, and $\mathrm C_k(f;s)=\omega_k(\vartheta(
\overline{O_kf})\tau_{s/a_k}O_kf)-|\omega_k(O_kf)|^2$.

\begin{proposition}[the spectral input]\label{prop:spectralinput}
Let $(\mathcal H,\Omega,H,\psi)$ be a continuum limit of $(\omega_k,O_k)$ and
$\Delta>0$. The following are equivalent.
\begin{enumerate}[label=\textup{(\roman*)},leftmargin=*]
\item $\sigma(H|_{\Omega^\perp})\subset[\Delta,\infty)$.
\item For every $f\in\mathcal T$ there is $C_f<\infty$ with $\lim_k\mathrm C_k(
f;s)\le C_fe^{-\Delta s}$ for every $s\in\mathbb D$.
\end{enumerate}
\end{proposition}

\begin{proof}
For $f\in\mathcal T$ let $\phi_f=\psi f-\ip\Omega{\psi f}\Omega$. Since $e^{-sH}
\Omega=\Omega$ and $e^{-sH}$ is self-adjoint, $\ip{\phi_f}{e^{-sH}\phi_f}=\ip{
\psi f}{e^{-sH}\psi f}-|\ip\Omega{\psi f}|^2$, so $\lim_k\mathrm C_k(f;s)$ exists
and equals $\ip{\phi_f}{e^{-sH}\phi_f}$. If (i) holds, the spectral theorem
gives (ii) with $C_f=\|\phi_f\|^2$. Conversely, let $\mu_f$ be the spectral
measure of $H$ at $\phi_f$. If $\mu_f([0,\Delta-\eta])=c>0$ for some $\eta>0$,
then $\ip{\phi_f}{e^{-sH}\phi_f}\ge ce^{-(\Delta-\eta)s}>C_fe^{-\Delta s}$ for
large $s\in\mathbb D$, which contradicts (ii); so the spectral projection $P$ of
$H$ for $[0,\Delta)$ satisfies $P\phi_f=0$. The $\phi_f$ are dense in
$\Omega^\perp$, so $P=0$ on $\Omega^\perp$.
\end{proof}

So the continuum gap needs exponential decay, at physical times and only
asymptotically in $k$, of the autocorrelations of the observables that generate
the limit; it does not need a lattice gap, since lattice states of physical
energy below $\Delta$ are allowed if they decouple from every $O_k(f)$.

\begin{lemma}\label{lem:limsupgap}
Let $(\mathcal H,\Omega,H,\psi)$ be a continuum limit of $(\omega_k,O_k)$ along
a subsequence $\Lambda$ and $\bar E=\limsup_{k\in\Lambda}m_k/a_k\in[0,\infty]$.
Then $\ip\phi{e^{-sH}\phi}\le e^{-\bar Es}\|\phi\|^2$ for all $\phi\in\Omega^\perp$
and $s\ge0$; so $\sigma(H|_{\Omega^\perp})\subset[\bar E,\infty)$, and $\Omega^\perp
=\{0\}$ if $\bar E=\infty$.
\end{lemma}

\begin{proof}
For $f\in\mathcal T$ and $k\in\Lambda$ put $O=O_kf$ and $\phi_k=\Psi_k(O)-
\omega_k(O)\Omega_k\perp\Omega_k$. Since $\mathbb T_k\ge0$ is self-adjoint and
fixes $\Omega_k$, it preserves $\Omega_k^\perp$, and for $s\in\mathbb D$ and
$t=s/a_k\in\Z$, $\mathrm C_k(f;s)=\ip{\phi_k}{\mathbb T_k^t\phi_k}\le\gamma_k^t
\|\phi_k\|^2=e^{-(m_k/a_k)s}\|\phi_k\|^2$. Along a further subsequence $m_k/a_k
\to\bar E$, $\mathrm C_k(f;s)\to\ip{\phi_f}{e^{-sH}\phi_f}$ and $\|\phi_k\|^2=
\mathrm C_k(f;0)\to\|\phi_f\|^2$, with $e^{-\infty s}=0$ for $s>0$. The $\phi_f$
are dense in $\Omega^\perp$, both sides are continuous in $\phi$, and $s\mapsto
\ip\phi{e^{-sH}\phi}$ is continuous, so the inequality holds on $\Omega^\perp$
for all $s\ge0$. The spectral statement follows as in
Proposition~\ref{prop:spectralinput}; if $\bar E=\infty$, then $\ip\phi{e^{-sH}
\phi}=0$ for $s>0$, so $\phi=0$ by strong continuity.
\end{proof}

\begin{theorem}[gap tuning]\label{thm:gaptuning}
Suppose $0<m_k<\infty$ for all $k$ and $m_k\to0$, and let $E>0$ and $a_0>0$.
For each $k$ let $j_k$ be the integer with $E\le m_k2^{j_k}/a_0<2E$, pass to a
subsequence on which $j_k$ is strictly increasing, and put $a_k=a_02^{-j_k}$.
\begin{enumerate}[label=\textup{(\alph*)},leftmargin=*]
\item Every continuum limit of $(\omega_k,O_k)$ along any subsequence, for any
$\mathcal T$ and $O_k$, has $\sigma(H|_{\Omega^\perp})\subset[E,\infty)$.
\item There are a subsequence, a space $\mathcal T$ with maps $O_k$, and a
nontrivial continuum limit along it with $\inf\sigma(H|_{\Omega^\perp})=\lim_k
m_k/a_k\in[E,2E]$.
\end{enumerate}
\end{theorem}

\begin{proof}
The integer $j_k$ exists and is unique, because the condition reads
\[
2^{j_k}\in[Ea_0/m_k,\,2Ea_0/m_k),
\]
and $j_k\to\infty$ since $m_k\to0$. Part (a) is
Lemma~\ref{lem:limsupgap}, since $m_k/a_k\ge E$.

(b) Pass to a subsequence with $m_k/a_k\to\bar E\in[E,2E]$ and fix $k$. As
$\gamma_k>0$, $\Omega_k^\perp\ne\{0\}$, and $\gamma_k$, the norm of the positive
operator $\mathbb T_k|_{\Omega_k^\perp}$, lies in its spectrum; so the spectral
projection $P_k$ of $\mathbb T_k|_{\Omega_k^\perp}$ for $[\gamma_ke^{-a_k/k},
\gamma_k]$ is nonzero. Choose a unit vector $v_k$ in its range. The vectors
$\Psi_k(F-\omega_k(F))$ are the projections of $\Psi_k(F)$ onto $\Omega_k^\perp$,
hence dense in it; choose $F'_k$ with $\|u'_k-v_k\|\le\frac1{10}$ for $u'_k=
\Psi_k(F'_k-\omega_k(F'_k))$, and put $F_k=(F'_k-\omega_k(F'_k))/\|u'_k\|$ and
$u_k=\Psi_k(F_k)$. Then $\|u_k\|=1$, $u_k\perp\Omega_k$ and $\|P_ku_k\|\ge(1-
\frac1{10})/\frac{11}{10}=\frac9{11}$. Let $c_k(t)=\ip{u_k}{\mathbb T_k^tu_k}$.
Since $\mathbb T_k^t\ge0$ commutes with $P_k$ and has norm $\gamma_k^t$ on
$\Omega_k^\perp$,
\begin{equation}\label{eq:twosided}
\tfrac{81}{121}e^{-(m_k+a_k/k)t}\le\ip{P_ku_k}{\mathbb T_k^tP_ku_k}\le c_k(t)\le
e^{-m_kt}.
\end{equation}
Let $\sigma_k$ be the spectral measure of $\mathbb T_k$ at $u_k$, a probability
measure on $[0,1]$, and $\rho_k$ its image under $x\mapsto-a_k^{-1}\log x$, a
probability measure on the compact space $[0,\infty]$, with $0\mapsto\infty$.
Then $c_k(s/a_k)=\int e^{-sx}\,d\rho_k(x)$ for $s\in\mathbb D$, with $e^{-s\infty}
=0$ for $s>0$. Pass to a subsequence with $\rho_k\to\rho$ weakly on $[0,\infty]$.
For $s>0$ the function $x\mapsto e^{-sx}$ is continuous on $[0,\infty]$, so with
$\rho'=\rho|_{[0,\infty)}$,
\[
B(s):=\lim_kc_k(s/a_k)=\int_{[0,\infty)}e^{-sx}\,d\rho'(x)\qquad(s\in\mathbb D,
\ s>0).
\]
Let $\mathcal T$ be the free complex vector space on the symbols $\mathbf1$ and
$j_s$, $s\in\mathbb D$, $s>0$; put $O_k(\mathbf1)=1$ and $O_k(j_s)=\tau_{s/a_k}
F_k$ for the $k$ with $s/a_k\in\Z$, extended linearly. Let $\mathcal H=\C\Omega
\oplus L^2([0,\infty),\rho')$ with $\Omega=(1,0)$, $H=0\oplus x$, $\psi(\mathbf1)
=\Omega$ and $\psi(j_s)=(0,e^{-sx})$. Since $\Psi_k(O_kj_s)=\mathbb T_k^{s/a_k}
u_k$, for $s,s'>0$ and $r$ in $\mathbb D$,
\[
\omega_k\bigl(\vartheta(\overline{O_kj_s})\tau_{r/a_k}O_kj_{s'}\bigr)=c_k\bigl((s
+s'+r)/a_k\bigr)\to B(s+s'+r)=\ip{\psi j_s}{e^{-rH}\psi j_{s'}};
\]
moreover $\omega_k(\tau_{r/a_k}O_kj_s)=\ip{\Omega_k}{\mathbb T_k^{(r+s)/a_k}u_k}=
\ip{\Omega_k}{u_k}=0=\ip\Omega{e^{-rH}\psi j_s}$, the same holds with the roles
exchanged, and $\omega_k(1)=1=\ip\Omega{e^{-rH}\Omega}$. By sesquilinearity the
identities of Definition~\ref{def:limit} hold on $\mathcal T$. The span of
$\{e^{-sx}:s\in\mathbb D,s>0\}$ is an algebra of functions in $C_0([0,\infty))$
that separates points and vanishes nowhere, hence dense in $C_0([0,\infty))$ by
the Stone--Weierstrass theorem and dense in $L^2(\rho')$; so $\psi$ has dense
range, and $H\ge0$ is self-adjoint with $H\Omega=0$.

By~\eqref{eq:twosided} with $t=s/a_k$, $\frac{81}{121}e^{-\bar Es}\le B(s)\le e^{
-\bar Es}$ for $s\in\mathbb D$, $s>0$. The lower bound gives $\rho'\ne0$, so
$\mathcal H\ne\C\Omega$. The upper bound gives $\rho'([0,\bar E-\eta])=0$ for
every $\eta>0$, and the lower bound gives $\rho'([\bar E,\bar E+\eta])>0$ for
every $\eta>0$. Since $\sigma(H|_{\Omega^\perp})$ is the support of $\rho'$, its
infimum is $\bar E$.
\end{proof}

Theorem~\ref{thm:gaptuning} uses no property of the gauge theory: it applies
equally to the free scalar field with lattice masses tending to zero. A
nontrivial continuum limit with a gap, in the sense of
Definition~\ref{def:limit}, therefore exists as soon as the lattice gaps
$m(\beta_k)$ are positive and tend to zero along a sequence $\beta_k\to\infty$,
and neither property is known for four-dimensional $\SU(N)$. The limit so
obtained is generated by one sequence of observables chosen through spectral
projections; nothing identifies it with Yang--Mills theory, since it carries no
Euclidean invariance, no locality and no Wilson loops. The physical content of
(E) is the convergence of fixed physical observables, smeared fields and Wilson
loops of fixed physical size, and by Lemma~\ref{lem:limsupgap} a gap coming
from the transfer matrices requires $m_k/a_k$ to stay bounded below along a
subsequence: along the asymptotic-freedom tuning, the lattice gap must decay at
exactly the asymptotic-freedom rate.
